\SourcePage{184}
\SGAsection{7}{纤维拓扑斯与纤维景}
\textbf{7.1. \Term{fibered-topos}{纤维拓扑斯}。}
\begin{SGAresult}{定义}{7.1.1}
设 $\mathcal U$ 为一个\Term{universe}{宇宙}。\Term{category}{范畴} $I$ 上的一个 $\mathcal U$-\Term{fibered-topos}{纤维拓扑斯}，是指一个 $I$ 上的\Term{fibered-category}{纤维范畴}（6.1）
\[p:F\longrightarrow I,\]
其\Term{fiber}{纤维}都是 $\mathcal U$-\Term{topos}{拓扑斯}，并且对于 $I$ 的\Term{morphism}{态射} $f:i\to j$，相应的\Term{inverse-image-functor}{逆像函子}\footnote{这里指\Term{fibered-structure}{纤维结构}的\Term{inverse-image}{逆像}（6.1.2）。}是某个\Term{topos-morphism}{拓扑斯态射} $f_\bullet:F_i\to F_j$（第 IV 讲 3.1.2）的\Term{inverse-image-functor}{逆像函子} $f^*:F_j\to F_i$。
\end{SGAresult}
\textbf{7.1.2.} \Term{equivalence}{等价}地，$I$ 上的一个 $\mathcal U$-\Term{fibered-topos}{纤维拓扑斯}，是指\Term{fiber}{纤维}为 $\mathcal U$-\Term{topos}{拓扑斯}、\Term{inverse-image-functor}{逆像函子}为\Term{exact-functor}{正合函子}且与\Term{inductive-limit}{归纳极限}交换的\Term{fibered-category}{纤维范畴} $p:F\to I$（第 IV 讲 1.6）。

\textbf{7.1.3.} 对 $I$ 的每个\Term{morphism}{态射} $f:i\to j$，选取一个\Term{inverse-image-functor}{逆像函子} $f^*:F_j\to F_i$。于是对于 $I$ 中每一对可复合\Term{morphism}{态射} $i\xrightarrow f j\xrightarrow g k$，有一个\Term{canonical}{典范}\Term{isomorphism}{同构} $c_{f,g}:f^*g^*\to(gf)^*$（6.1.2.1），而这些 $c_{f,g}$ 满足\Term{cocycle}{上循环}性质（6.1.2.2）。此外，对每个 $f:i\to j$，为 $f^*:F_j\to F_i$ 选取一个\Term{right-adjoint}{右伴随函子} $f_*:F_i\to F_j$。作为 $f^*$ 的\Term{right-adjoint}{右伴随函子}这一性质将 $f_*$ 确定到\Term{canonical}{典范}\Term{isomorphism}{同构}。由\Term{adjoint-functor}{伴随函子}的一般结果，得到\Term{canonical}{典范}\Term{isomorphism}{同构}
\[c'_{g,f}:g_*f_*\longrightarrow(gf)_*.\tag{7.1.3.1}\]
\SourcePage{185}
它们满足与 6.2.2.2 类似的\Term{cocycle}{上循环}性质，其明确表述留给读者。应用 6.1.3 的结果，便得到一个 $I^\circ$ 上的\Term{fibered-category}{纤维范畴}
\[p':F'\longrightarrow I^\circ.\tag{7.1.3.2}\]
容易验证，\Term{fibered-category}{纤维范畴} $p':F'\to I^\circ$ 不依赖于构造时的各项选择，精确到唯一的 $I^\circ$-\Term{isomorphism}{同构}\footnote{执笔者为这一构造缺乏内在性而致歉。}。它在 $I^\circ$ 的每个\Term{object}{对象} $i$ 处的\Term{fiber}{纤维}\Term{canonical}{典范}\Term{isomorphism}{同构}于 $\mathcal U$-\Term{topos}{拓扑斯} $F_i$，并与后者等同。\Term{fibered-category}{纤维范畴} $p':F'\to I^\circ$ 的\Term{inverse-image-functor}{逆像函子}是\Term{topos-morphism}{拓扑斯态射}的\Term{direct-image-functor}{正像函子}。

\textbf{7.1.4.} 设 $p:F\to I$ 为一个 $\mathcal U$-\Term{fibered-topos}{纤维拓扑斯}。对于每个 $f:i\to j$，选取一个\Term{topos-morphism}{拓扑斯态射} $f_\bullet:F_i\to F_j$（第 IV 讲 3.1.2），使其\Term{inverse-image-functor}{逆像函子} $f^*$ 是\Term{fibered-structure}{纤维结构}的\Term{inverse-image}{逆像}（6.1）。于是得到联系这些\Term{topos-morphism}{拓扑斯态射}的\Term{cocycle}{上循环} $c_{f,g}$（7.1.3.1 和 7.1.3.2）。\Term{morphism}{态射} $f_\bullet$，$f\in\operatorname{Fl}(I)$，及\Term{cocycle}{上循环} $c_{f,g}$ 的全体称为\Term{fibered-topos}{纤维拓扑斯} $F\to I$ 的一个\Term{bisplitting}{双分裂}（biscindage）。选定一个\Term{bisplitting}{双分裂}后，便可将一个从\Term{category}{范畴} $I$ 到 $\mathcal U$-\Term{topos}{拓扑斯}的 $2$-\Term{category}{范畴}的\Term{pseudofunctor}{伪函子} $i\mapsto F_i$ 与 $F\to I$ 相联系。反过来，可本质唯一地（参见 [5]）将一个带有\Term{bisplitting}{双分裂}的 $\mathcal U$-\Term{fibered-topos}{纤维拓扑斯} $F\to I$ 与之相联系，使对应的\Term{pseudofunctor}{伪函子}\Term{canonical}{典范}\Term{isomorphism}{同构}于 $i\mapsto F_i$。实际中通常将 $\mathcal U$-\Term{fibered-topos}{纤维拓扑斯} $F\to I$ 记为 $(F_i)_{i\in I}$，并默认已经选定一个\Term{bisplitting}{双分裂}。但须注意，对\Term{fibered-topos}{纤维拓扑斯}所作的构造，如相应的\Term{total-topos}{全拓扑斯}（7.4）和\Term{projective-limit}{投射极限}（8.1），只依赖于\Term{fibered-topos}{纤维拓扑斯}本身，而不依赖于可能选取的\Term{bisplitting}{双分裂}。

\begin{SGAresult}{定义}{7.1.5}
设 $p:F\to I$ 和 $q:G\to I$ 为 $I$ 上的两个 $\mathcal U$-\Term{fibered-topos}{纤维拓扑斯}。从 $(F,p)$ 到 $(G,q)$ 的一个\Term{fibered-topos}{纤维拓扑斯}\Term{morphism}{态射} $m$，由一个 $I$-\Term{functor}{函子} $m^*:G\to F$（6.1.4）及如下数据组成：对 $I$ 的每个\Term{object}{对象} $i$，给定 $m^*$ 在 $i$ 处\Term{fiber}{纤维}上的\Term{restriction}{限制} $m_i^*:G_i\to F_i$ 的一个\Term{right-adjoint}{右伴随函子} $m_{i*}:F_i\to G_i$。这些数据还须满足：对 $I$ 的每个\Term{object}{对象} $i$，\Term{adjoint-functor}{伴随函子}对 $(m_i^*,m_{i*})$ 是从\Term{topos}{拓扑斯} $F_i$ 到\Term{topos}{拓扑斯} $G_i$ 的一个\Term{morphism}{态射}。
\end{SGAresult}
\textbf{7.1.6.} 设 $m:(F,p)\to(G,q)$ 是 $I$ 上的一个 $\mathcal U$-\Term{fibered-topos}{纤维拓扑斯}\Term{morphism}{态射}，并设 $p':F'\to I^\circ$ 和 $q':G'\to I^\circ$ 分别由 $(F,p)$ 和 $(G,q)$ 按 7.1.3 的方法得到。选取 $F$ 和 $G$ 的\Term{bisplitting}{双分裂}（7.1.4），两者都记为 $(f:i\to j)\mapsto f_\bullet=(f^*,f_*)$，分别是 $F_i\to F_j$ 和 $G_i\to G_j$ 的\Term{topos-morphism}{拓扑斯态射}。因为 $f^*:G\to F$ 是一个 $I$-\Term{functor}{函子}，对于每个 $f:i\to j$，有一个\Term{functor}{函子}间\Term{morphism}{态射}（\Term{transition-morphism}{转移态射}）
\[b_f:f^*m_j^*\longrightarrow m_i^*f^*.\tag{7.1.6.1}\]
\SourcePage{186}
对于每对可复合\Term{morphism}{态射} $i\xrightarrow f j\xrightarrow g k$，有
\[m_i^*(c_{f,g})b_f^*f^*(b_g)=b_{gf}c_{f,g}.\tag{7.1.6.2}\]
转到\Term{adjoint-functor}{伴随函子}，得到\Term{morphism}{态射}
\[b'_f:f_*m_{i*}\longrightarrow m_{j*}f_*\tag{7.1.6.3}\]
满足与 (7.1.6.2) 类似的关系。由于这些 $f_*$ 是\Term{fibered-category}{纤维范畴} $(F',p')$ 和 $(G',q')$ 的\Term{inverse-image-functor}{逆像函子}（7.1.3），可以构造一个 $I^\circ$-\Term{functor}{函子}
\[m_*:F'\longrightarrow G',\tag{7.1.6.4}\]
其在\Term{fiber}{纤维}上的\Term{restriction}{限制}是 $m_{i*}$，而\Term{transition-isomorphism}{转移同构}是 $b'_f$。稍加耐心便可验证，这样构造的\Term{functor}{函子} $m_*$ 不依赖于\Term{topos-morphism}{拓扑斯态射} $f_\bullet$ 的选择\footnote{参见原书第 179 页的脚注（本块 PDF 第 185 页）。}。

\textbf{7.1.7.} 以 $\operatorname{Homtop}_I(E,G)$ 表示\Term{fibered-topos}{纤维拓扑斯} $(F,p)$ 到 $(G,q)$ 的\Term{morphism}{态射}\Term{set}{集合}。它是一个\Term{category}{范畴} $\mathscr{H}\!omtop_I(F,G)$ 的\Term{object}{对象}集：若 $m,n$ 为两个\Term{fibered-topos}{纤维拓扑斯}\Term{morphism}{态射}，则从 $m$ 到 $n$ 的一个\Term{morphism}{态射}是从 $n^*$ 到 $m^*$ 的一个 $I$-\Term{functor}{函子}\Term{morphism}{态射}。通过\Term{adjunction}{伴随}，还得到从 $m_*$ 到 $n_*$ 的一个 $I^\circ$-\Term{morphism}{态射}（7.1.6）。因此最终得到两个\Term{functor}{函子}
\[\mathscr{H}\!omtop_I(F,G)\longrightarrow\mathscr{H}\!om_I(G,F)^\circ,\tag{7.1.7.1}\]
\[\mathscr{H}\!omtop_I(F,G)\longrightarrow\mathscr{H}\!om_{I^\circ}(F',G'),\tag{7.1.7.2}\]
它们将每个\Term{morphism}{态射} $(m^*,(m_{i*})_{i\in\operatorname{Ob}I})$ 分别映到\Term{functor}{函子} $m^*\in\operatorname{Hom}_I(G,F)$ 和按 7.1.6 将 $m_{i*}$ \Term{gluing}{粘合}而成的\Term{functor}{函子} $m_*\in\operatorname{Hom}_{I^\circ}(F',G')$。\Term{functor}{函子} $m^*:G\to F$（7.1.7.1）称为\Term{fibered-topos}{纤维拓扑斯}\Term{morphism}{态射} $m:F\to G$ 的\Term{inverse-image-functor}{逆像函子}；\Term{functor}{函子} $m_*:F'\to G'$（7.1.7.2）称为其\Term{direct-image-functor}{正像函子}。定义立即表明，\Term{functor}{函子} $m\mapsto m^*$ 和 $m\mapsto m_*$ 均\Term{fully-faithful}{完全忠实}，其\Term{essential-image}{本质像}一方面是逐\Term{fiber}{纤维}\Term{exact-functor}{正合}并与\Term{inductive-limit}{归纳极限}交换的 $G\to F$ 的全部 $I$-\Term{functor}{函子}，另一方面是逐\Term{fiber}{纤维}\Term{exact-functor}{正合}并与\Term{inductive-limit}{归纳极限}交换的 $F'\to G'$ 的全部 $I^\circ$-\Term{functor}{函子}，另一方面是逐\Term{fiber}{纤维}具有\Term{exact-functor}{正合}\Term{left-adjoint}{左伴随函子}的 $F'\to G'$ 的全部 $I^\circ$-\Term{functor}{函子}。这说明以\Term{direct-image-functor}{正像函子}或\Term{inverse-image-functor}{逆像函子}来定义\Term{fibered-topos}{纤维拓扑斯}\Term{morphism}{态射}的用语是合理的。事实上，这样的\Term{morphism}{态射}只是确定到唯一\Term{isomorphism}{同构}。

以 $\mathscr{H}\!omtop_{\mathrm{Cart}/I}(F,G)$ 表示由满足 $m_*$（\Term{equivalence}{等价}地 $m^*$）为\Term{cartesian-functor}{笛卡尔函子}（6.1.4）的\Term{morphism}{态射} $m$ 生成的 $\mathscr{H}\!omtop_I(F,G)$ 的\Term{full-subcategory}{满子范畴}。
\SourcePage{187}
这种\Term{morphism}{态射}称为\Term{cartesian-morphism}{笛卡尔态射}。于是由 (7.1.7.1)、(7.1.7.2) 诱导\Term{fully-faithful}{完全忠实}\Term{functor}{函子}
\[\mathscr{H}\!omtop_{\mathrm{Cart}/I}(F,G)\longrightarrow\mathscr{H}\!om_{\mathrm{cart}/I}(G,F)^\circ,\tag{7.1.7.3}\]
\[\mathscr{H}\!omtop_{\mathrm{Cart}/I}(F,G)\longrightarrow\mathscr{H}\!om_{\mathrm{cart}/I^\circ}(F',G').\tag{7.1.7.4}\]
\textbf{7.1.8.} 设 $p:F\to I$、$q:G\to I$、$r:H\to I$ 为 $I$ 上的三个 $\mathcal U$-\Term{fibered-topos}{纤维拓扑斯}，$m:F\to G$、$n:G\to H$ 为两个\Term{morphism}{态射}。如同第 IV 讲 3.3，定义其复合，记为 $nm:F\to H$。若 $\mathcal V$ 是满足 $\mathcal U\in\mathcal V$ 的\Term{universe}{宇宙}，则可定义\Term{category}{范畴} $(\mathcal V\text{-}\mathcal U\text{-}\mathrm{Top}/I)$，其\Term{object}{对象}是\Term{fiber}{纤维}\Term{category}{范畴}属于 $\mathcal V$ 的 $I$ 上 $\mathcal U$-\Term{fibered-topos}{纤维拓扑斯}，其\Term{morphism}{态射}是 $I$ 上的 $\mathcal U$-\Term{fibered-topos}{纤维拓扑斯}\Term{morphism}{态射}。此外，如同第 IV 讲 3.3.2，复合映射
\[\operatorname{Homtop}_{/I}(F,G)\times\operatorname{Homtop}_{/I}(G,F)\longrightarrow\operatorname{Homtop}_{/I}(F,H)\]
是一个“\Term{morphism}{态射}复合\Term{functor}{函子}”
\[\mathscr{H}\!omtop_{/I}(F,G)\,\mathscr{H}\!omtop_{/I}(G,H)\longrightarrow\mathscr{H}\!omtop_{/I}(F,H)\]
在\Term{object}{对象}上诱导的映射；第 IV 讲 3.3、3.4 的讨论原封不动地推广到 $\mathcal U$-\Term{fibered-topos}{纤维拓扑斯}。读者也许已经注意到，当 $I$ 是一个点\Term{category}{范畴}（一个\Term{object}{对象}，一个恒等\Term{morphism}{态射}）时，$I$ 上的一个 $\mathcal U$-\Term{fibered-topos}{纤维拓扑斯}“无非”是一个 $\mathcal U$-\Term{topos}{拓扑斯}\Term{morphism}{态射}，而 $\mathcal U$-\Term{fibered-topos}{纤维拓扑斯}的一个\Term{morphism}{态射}“无非”是一个 $\mathcal U$-\Term{topos}{拓扑斯}\Term{morphism}{态射}，上述构造便归结为第 IV 讲 3 的构造。

\textbf{7.1.9.} 设 $p:F\to I$ 是一个 $\mathcal U$-\Term{fibered-topos}{纤维拓扑斯}，$\phi:J\to I$ 是一个\Term{functor}{函子}。由 $(F,p)$ 经\Term{base-change}{基变换} $\phi$ 得到的\Term{fibered-category}{纤维范畴} $p_J:F_J\to J$ 是 $J$ 上的一个 $\mathcal U$-\Term{fibered-topos}{纤维拓扑斯}。设 $p':F'\to I^\circ$ 和 $(p_J)':(F_J)'\to J^\circ$ 分别按 7.1.3 从 $(F,p)$ 和 $(F_J,p_J)$ 构造。立即可验证，$((F_J)',(p_J)')$ \Term{canonical}{典范}\Term{isomorphism}{同构}于由 $p':F'\to I^\circ$ 经\Term{base-change}{基变换} $\phi:J^\circ\to I^\circ$ 得到的 $(p')_J:(F')_J\to J^\circ$。以下将这两个 $J^\circ$ 上的\Term{fibered-category}{纤维范畴}等同，记为 $p'_J:F'_J\to J^\circ$。\Term{base-change}{基变换}关于其自变量具有\Term{functorial}{函子性}，甚至具有 $2$-\Term{functorial}{函子性}：对于 $I$ 上每对 $\mathcal U$-\Term{fibered-topos}{纤维拓扑斯} $p:F\to I$、$G\to I$，有两个\Term{category}{范畴}与\Term{functor}{函子}的\Term{commutative-diagram}{交换图}
\[\begin{tikzcd}
\mathscr{H}\!omtop_I(F,G)\arrow[r]\arrow[d]&\mathscr{H}\!om_I(G,F)^\circ\arrow[d]\\
\mathscr{H}\!omtop_J(F_J,G_J)\arrow[r]&\mathscr{H}\!om_J(G_J,F_J)^\circ
\end{tikzcd}\tag{7.1.9.1}\]
\SourcePage{188}
\[\begin{tikzcd}
\mathscr{H}\!omtop_I(F,G)\arrow[r]\arrow[d]&\mathscr{H}\!om_{I^\circ}(F',G')\arrow[d]\\
\mathscr{H}\!omtop_J(F_J,G_J)\arrow[r]&\mathscr{H}\!om_{J^\circ}(F'_J,G'_J)
\end{tikzcd}\tag{7.1.9.2}\]
其中竖直\Term{functor}{函子}是\Term{base-change}{基变换}\Term{functor}{函子}，水平\Term{functor}{函子}在 (7.1.9.1) 中是“\Term{morphism}{态射} $\mapsto$ \Term{inverse-image}{逆像}”\Term{functor}{函子}（7.1.7.1），在 (7.1.9.2) 中是“\Term{morphism}{态射} $\mapsto$ \Term{direct-image}{正像}”\Term{functor}{函子}（7.1.7.2）。

\textbf{7.2. \Term{fibered-site}{纤维景}。}

\textbf{7.2.1.} \Term{category}{范畴} $I$ 上的一个 $\mathcal U$-\Term{fibered-site}{纤维景} $C$，是一个\Term{fibered-category}{纤维范畴} $p:C\to I$，其\Term{fiber}{纤维}各配备一个\Term{topology}{拓扑}而成为 $\mathcal U$-\Term{site}{景}（第 IV 讲 3.0.2），并且对于 $I$ 中每个\Term{morphism}{态射} $f:i\to j$，\Term{inverse-image-functor}{逆像函子} $f^*:C_j\to C_i$ 是从\Term{site}{景} $C_i$ 到\Term{site}{景} $C_j$ 的一个\Term{site-morphism}{景态射}（第 IV 讲 4.9.1）。

\textbf{7.2.2.} 设 $p:C\to I$、$q:D\to I$ 为 $I$ 上的两个 $\mathcal U$-\Term{fibered-site}{纤维景}。从 $C$ 到 $D$ 的一个 $\mathcal U$-\Term{fibered-site}{纤维景}\Term{morphism}{态射}，是指一个 $I$-\Term{functor}{函子}（6.0）$m:D\to C$，使每个\Term{object}{对象} $i$ 上的 $m_i:D_i\to C_i$ 都是从\Term{site}{景} $C_i$ 到\Term{site}{景} $D_i$ 的一个\Term{site-morphism}{景态射}。这些\Term{morphism}{态射}的\Term{set}{集合}记为 $\operatorname{Morsite}_I(C,D)$。从\Term{fibered-site}{纤维景} $D$ 到\Term{fibered-site}{纤维景} $C$ 的一个\Term{continuous-functor}{连续函子}，是指逐\Term{fiber}{纤维}诱导\Term{continuous-functor}{连续函子}（第 III 讲 1）的 $I$-\Term{functor}{函子} $D\to C$。其\Term{set}{集合}记为 $\operatorname{Cont}_I(D,C)$。以 $\mathscr{M}orsite_I(C,D)$ 表示由\Term{object}{对象}集 $\operatorname{Morsite}_I(C,D)\subset\operatorname{Hom}_I(D,C)$ 定义的 $\mathscr{H}\!om_I(D,C)^\circ$ 的\Term{full-subcategory}{满子范畴}。类似地，以 $\mathscr{C}ont_I(D,C)$ 表示由\Term{object}{对象}\Term{subset}{子集} $\operatorname{Cont}_I(D,C)\subset\operatorname{Hom}_I(D,C)$ 定义的 $\mathscr{H}\!om_I(D,C)$ 的\Term{full-subcategory}{满子范畴}。因而有两个\Term{fully-faithful}{完全忠实}且在\Term{object}{对象}上\Term{injectivity}{单射}的\Term{functor}{函子}
\[\mathscr{M}orsite_I(C,D)\hookrightarrow\mathscr{C}ont_I(D,C)^\circ\hookrightarrow\mathscr{H}\!om_I(D,C)^\circ.\tag{7.2.2.1}\]
如 7.1.7 一样定义从\Term{fibered-site}{纤维景} $C$ 到 $D$ 的\Term{cartesian-morphism}{笛卡尔态射}。有如下\Term{full-subcategory}{满子范畴}图：
\[\begin{tikzcd}
\mathscr{M}orsite_{\mathrm{Cart}/I}(C,D)\arrow[r,hook]\arrow[d,hook]&\mathscr{H}\!om_{\mathrm{Cart}/I}(D,C)^\circ\arrow[d,hook]\\
\mathscr{M}orsite_I(C,D)\arrow[r,hook]&\mathscr{H}\!om_I(D,C)^\circ.
\end{tikzcd}\tag{7.2.2.2}\]
\textbf{7.2.3.} $\mathcal U$-\Term{fibered-site}{纤维景}的\Term{morphism}{态射}可以复合。这给出\Term{category}{范畴} $(\mathcal V\text{-}\mathcal U\text{-}\mathrm{Site}/I)$，其\Term{object}{对象}是\Term{fiber}{纤维}属于\Term{universe}{宇宙} $\mathcal V$（且 $\mathcal U\in\mathcal V$）的 $\mathcal U$-\Term{site}{景}。
\SourcePage{189}
这个 $\mathcal V$-$\mathcal U$-\Term{site}{景}\Term{category}{范畴}还具有 $2$-\Term{category}{范畴}结构；特别地，$\mathcal U$-\Term{fibered-site}{纤维景}\Term{morphism}{态射}的复合关于\Term{morphism}{态射}之间的\Term{morphism}{态射}具有\Term{functorial}{函子性}。对每个\Term{functor}{函子} $\phi:J\to I$ 和每个\Term{fibered-site}{纤维景} $p:C\to I$，从 $(C,p)$ 经\Term{base-change}{基变换} $\phi$ 得到的\Term{category}{范畴} $p_J:C_J\to J$ \Term{canonical}{典范}地具有 $J$ 上的 $\mathcal U$-\Term{fibered-site}{纤维景}结构。\Term{base-change}{基变换}运算具有 $2$-\Term{functorial}{函子性}。

\textbf{7.2.4.} 设 $p:C\to I$ 是一个\Term{fibered-category}{纤维范畴}，其\Term{fiber}{纤维}各赋予\Term{topology}{拓扑}而成为 $\mathcal U$-\Term{site}{景}。假设\Term{fibered-category}{纤维范畴}中的\Term{projective-limit}{投射极限}\Term{representable}{可表}。为了使 $p:C\to I$ 成为 $\mathcal U$-\Term{fibered-site}{纤维景}，只须\Term{inverse-image-functor}{逆像函子}\Term{left-exact}{左正合}且将\Term{covering-family}{覆盖族}映成\Term{covering-family}{覆盖族}；当各\Term{fiber}{纤维}的\Term{topology}{拓扑}不细于\Term{canonical-topology}{典范拓扑}时，这个条件也是充分的（第 IV 讲 4.9.2）。本研讨班所用的\Term{fibered-site}{纤维景}都具有上述形式。类似地，设 $p:C\to I$、$q:D\to I$ 为两个\Term{fibered-site}{纤维景}，其\Term{fiber}{纤维}\Term{category}{范畴}中的\Term{finite}{有限}\Term{projective-limit}{投射极限}\Term{representable}{可表}。若\Term{cartesian-functor}{笛卡尔函子} $m:D\to C$ 在\Term{fiber}{纤维}\Term{category}{范畴}上诱导\Term{left-exact}{左正合}\Term{functor}{函子}且将\Term{fiber}{纤维}\Term{category}{范畴}的\Term{covering-family}{覆盖族}映成\Term{covering-family}{覆盖族}，则它是从 $C$ 到 $D$ 的\Term{fibered-site}{纤维景}\Term{morphism}{态射}；若 $C$ 各\Term{fiber}{纤维}的\Term{topology}{拓扑}不细于\Term{canonical-topology}{典范拓扑}，逆命题亦成立（第 IV 讲 4.9.2）。本研讨班所用的\Term{fibered-site}{纤维景}\Term{morphism}{态射}都具有上述形式。

\textbf{7.2.5.} 在各\Term{fiber}{纤维}上赋予\Term{canonical-topology}{典范拓扑}时，$\mathcal U$-\Term{fibered-topos}{纤维拓扑斯} $p:F\to I$ 是一个 $\mathcal U$-\Term{fibered-site}{纤维景}；以下总采用此约定。设 $m:F\to G$ 是 $\mathcal U$-\Term{fibered-topos}{纤维拓扑斯}的一个\Term{morphism}{态射}。\Term{inverse-image-functor}{逆像函子} $m^*:G\to F$（7.1.7）是从 $\mathcal U$-\Term{fibered-site}{纤维景} $F$ 到 $\mathcal U$-\Term{fibered-site}{纤维景} $G$ 的一个\Term{morphism}{态射}。于是得到一个\Term{functor}{函子} $(\mathcal V\text{-}\mathcal U\text{-}\mathrm{Top}/I)\to(\mathcal V\text{-}\mathcal U\text{-}\mathrm{Site}/I)$；它一般不忠实。但实际上，这个\Term{functor}{函子}自然地延伸为一个 $2$-\Term{functor}{函子}，对于 $I$ 上的两个 $\mathcal U$-\Term{fibered-topos}{纤维拓扑斯} $F,G$，它诱导\Term{functor}{函子}
\[\mathscr{H}\!omtop_I(F,G)\longrightarrow\mathscr{M}orsite_I(F,G),\tag{7.2.5.1}\]
由定义立即看出，这是一个\Term{equivalence-of-categories}{范畴等价}。此外，将 (7.2.5.1) 与\Term{canonical}{典范}包含 $\mathscr{M}orsite_I(F,G)\hookrightarrow\mathscr{H}\!om_I(G,F)^\circ$ 复合，所得正是 (7.1.7.1)，即将\Term{fibered-topos}{纤维拓扑斯}\Term{morphism}{态射}映到其\Term{inverse-image}{逆像}的\Term{functor}{函子}。

\textbf{7.2.6.} 设 $p:C\to I$ 是一个 $\mathcal U$-\Term{fibered-site}{纤维景}。对于 $I$ 中每个\Term{morphism}{态射} $f:i\to j$，选取一个\Term{base-change}{基变换}\Term{functor}{函子} $f^*:C_i\to C_i$。转到 $\mathcal U$-\Term{sheaf}{层}\Term{category}{范畴}，$f^*$ 延伸为一个\Term{functor}{函子} $\operatorname{Top}(f)^*:C_j^\sim\to C_i^\sim$，其中 $\operatorname{Top}(f):C_i^\sim\to C_j^\sim$ 为\Term{topos-morphism}{拓扑斯态射}（第 IV 讲 4.9.1.1），使下图精确到\Term{isomorphism}{同构}交换：
\SourcePage{190}
\[\begin{tikzcd}
C_j\arrow[r,"f^*"]\arrow[d,"\varepsilon_j"']&C_i\arrow[d,"\varepsilon_i"]\\
C_j^\sim\arrow[r,"\operatorname{Top}(f)^*"']&C_i^\sim.
\end{tikzcd}\tag{7.2.6.1}\]
令
\[c_{f,g}:f^*g^*\longrightarrow(gf)^*\tag{7.2.6.2}\]
为\Term{canonical}{典范}\Term{isomorphism}{同构}（6.1.2.1），令
\[b_f:\operatorname{Top}(f)^*\varepsilon_j\longrightarrow\varepsilon_i f^*\tag{7.2.6.3}\]
为使 (7.2.6.1) “交换”的\Term{isomorphism}{同构}。利用\Term{functor}{函子} $\operatorname{Top}(f)^*$ 与\Term{inductive-limit}{归纳极限}交换这一事实，可立即验证：对每对可复合\Term{morphism}{态射} $i\xrightarrow f j\xrightarrow g k$，存在唯一\Term{isomorphism}{同构}
\[\operatorname{Top}(c_{f,g}):\operatorname{Top}(f)^*\operatorname{Top}(g)^*\longrightarrow\operatorname{Top}(gf)^*\tag{7.2.6.4}\]
使得
\[\varepsilon_i(c_{f,g})b_f\operatorname{Top}(f)^*(b_g)=b_{gf}\operatorname{Top}(c_{f,g}).\tag{7.2.6.5}\]
而且 (7.2.6.4) 中的\Term{isomorphism}{同构}自动满足 6.2.2.2 的\Term{cocycle}{上循环}条件。因此\Term{constructible}{可构造}一个\Term{fibered-category}{纤维范畴}
\[p^\sim:C^{\sim/I}\longrightarrow I,\tag{7.2.6.6}\]
其\Term{fiber}{纤维}\Term{canonical}{典范}\Term{isomorphism}{同构}于 $\mathcal U$-\Term{topos}{拓扑斯} $C_i^\sim$（6.1.3），并构造一个\Term{cartesian-functor}{笛卡尔函子}
\[\varepsilon_{C/I}:C\longrightarrow C^{\sim/I},\tag{7.2.6.7}\]
在\Term{fiber}{纤维}上诱导 $\varepsilon_i:C_i\to C_i^\sim$。定义立即表明，$p^\sim:C^{\sim/I}\to I$ 是 $I$ 上的一个 $\mathcal U$-\Term{fibered-topos}{纤维拓扑斯}，$\varepsilon_{C/I}$ 是从 $C^{\sim/I}$ 到 $C$ 的一个\Term{cartesian-morphism}{笛卡尔} $\mathcal U$-\Term{fibered-site}{纤维景}\Term{morphism}{态射}，而且 $C^{\sim/I}\to I$ 和 $\varepsilon_{C/I}$ 不依赖于构造中各项选择（\Term{base-change}{基变换}\Term{functor}{函子} $f^*$ 的选择及其延伸 $\operatorname{Top}(f)^*$ 的选择），精确到\Term{canonical}{典范}\Term{isomorphism}{同构}。$\mathcal U$-\Term{fibered-topos}{纤维拓扑斯} $C^{\sim/I}$ 称为与此 $\mathcal U$-\Term{fibered-site}{纤维景}相伴的 $\mathcal U$-\Term{fibered-topos}{纤维拓扑斯}。
\SourcePage{191}
\textbf{7.2.6.8.} 可验证，与\Term{fibered-site}{纤维景}相伴的 $\mathcal U$-\Term{fibered-topos}{纤维拓扑斯}的形成与底\Term{category}{范畴}的\Term{base-change}{基变换}运算“交换”。

\begin{SGAresult}{命题}{7.2.7}
设 $E$ 为 $I$ 上的一个 $\mathcal U$-\Term{fibered-topos}{纤维拓扑斯}，$C$ 为 $I$ 上的一个 $\mathcal U$-\Term{fibered-site}{纤维景}。将每个\Term{fibered-topos}{纤维拓扑斯}\Term{morphism}{态射} $m:E\to C^{\sim/I}$ 映到其\Term{inverse-image-functor}{逆像函子} $m^*:C^{\sim/I}\to E$ 与 $\varepsilon_{C/I}$ 的复合的\Term{functor}{函子} $m\mapsto m^*\varepsilon_{C/I}$，诱导一个保持\Term{cartesian-morphism}{笛卡尔对象}的\Term{equivalence-of-categories}{范畴等价}
\[\mathscr{H}\!omtop_I(E,C^{\sim/I})\xrightarrow\sim\mathscr{M}orsite_I(E,C).\]
若 $C$ 的\Term{fiber}{纤维}\Term{category}{范畴}中\Term{finite}{有限}\Term{projective-limit}{投射极限}\Term{representable}{可表}，则对应的\Term{fully-faithful}{完全忠实}\Term{functor}{函子}
\[\mathscr{H}\!omtop_I(E,C^{\sim/I})\longrightarrow\mathscr{H}\!om_I(C,E)^\circ\]
的\Term{essential-image}{本质像}为所有逐\Term{fiber}{纤维}\Term{left-exact}{左正合}、将\Term{fiber}{纤维}\Term{category}{范畴}的\Term{covering-family}{覆盖族}映成相应\Term{fiber}{纤维}\Term{category}{范畴}的\Term{epimorphic-family}{满射族}的 $I$-\Term{functor}{函子} $g:C\to E$。
\end{SGAresult}
\begin{proof}
此命题只是将第 IV 讲命题 4.9.4 推广到\Term{fiber}{纤维}情形；该处的证明可立即推广。
\end{proof}
\textbf{7.3. 例子。}
\begin{SGAresult}{例}{7.3.1}
（一个\Term{topos}{拓扑斯}的\Term{morphism}{态射}\Term{fibered-topos}{纤维拓扑斯}。）设 $E$ 为一个 $\mathcal U$-\Term{topos}{拓扑斯}，$\operatorname{Fl}(E)$ 为 $E$ 的\Term{morphism}{态射}\Term{category}{范畴}，$p:\operatorname{Fl}(E)\to E$ 为将\Term{morphism}{态射}映到其终点的\Term{functor}{函子}。\Term{category}{范畴} $(\operatorname{Fl}(E),p)$ 是 $E$ 上的一个 $\mathcal U$-\Term{fibered-topos}{纤维拓扑斯}（7.1.1 和第 IV 讲 5.1）。在 $E$ 的每个\Term{object}{对象} $X$ 处的\Term{fiber}{纤维}是\Term{topos}{拓扑斯} $E_{/X}$。\Term{morphism}{态射} $f:X\to Y$ 的\Term{inverse-image-functor}{逆像函子} $f^*:E_{/Y}\to E_{/X}$ 是\Term{fibre-product}{纤维积}\Term{functor}{函子}。所以对于每个\Term{functor}{函子} $\phi:I\to E$，可由\Term{base-change}{基变换} $\phi$ 得到 $I$ 上的一个 $\mathcal U$-\Term{fibered-topos}{纤维拓扑斯}
\[p_I:\operatorname{Fl}(E)_I\longrightarrow I,\tag{7.3.1.1}\]
其在 $I$ 的每个\Term{object}{对象} $i$ 处的\Term{fiber}{纤维}为 $E_{/\phi(i)}$（7.1.9）。若 $\mathcal V$ 是以 $E$ 为元素的\Term{universe}{宇宙}，则采用 7.1.8 的记号有一个映射
\[\operatorname{Hom}(I,E)\longrightarrow\operatorname{ob}(\mathcal V\text{-}\mathcal U\text{-}\mathrm{top}/I),\tag{7.3.1.2}\]
将 $\phi\in\operatorname{Hom}(I,E)$ 映到 $(\operatorname{Fl}(E),p)$ 经\Term{base-change}{基变换} $\phi$ 得到的 $I$ 上 $\mathcal U$-\Term{fibered-topos}{纤维拓扑斯}。设 $\phi_1,\phi_2:I\to E$ 为两个\Term{functor}{函子}，$m:\phi_1\to\phi_2$ 为\Term{functor}{函子}间的一个\Term{morphism}{态射}。以 $p_1:\operatorname{Fl}(E)_1\to I$ 和 $p_2:\operatorname{Fl}(E)_2\to I$ 分别表示由 $p$ 经\Term{base-change}{基变换} $\phi_1$、$\phi_2$ 得到的 $\mathcal U$-\Term{fibered-topos}{纤维拓扑斯}。对每个\Term{object}{对象} $i$，选取一个\Term{topos-morphism}{拓扑斯态射}（第 IV 讲 5.5.2）
\[\operatorname{loc}(m(i)):E_{/\phi_1(i)}\longrightarrow E_{/\phi_2(i)}.\tag{7.3.1.3}\]
\SourcePage{192}
容易验证，存在唯一\footnote{唯一性来自采用第 IV 讲 5.5.3 的记号时所要求的公式 $(gf)_!=g_!f_!$。}的 $I$ 上\Term{fibered-topos}{纤维拓扑斯}\Term{cartesian-morphism}{笛卡尔态射}
\[\operatorname{loc}(m):\operatorname{Fl}(E)_1\longrightarrow\operatorname{Fl}(E)_2,\tag{7.3.1.4}\]
在各\Term{fibered-topos}{纤维拓扑斯}上诱导\Term{morphism}{态射} $\operatorname{loc}(m(i))$。$\operatorname{loc}(m)$ 由\Term{functor}{函子}\Term{morphism}{态射} $m$ 确定到\Term{canonical}{典范}\Term{isomorphism}{同构}。若 $n:\phi_2\to\phi_3$ 为\Term{functor}{函子}\Term{morphism}{态射}，则有\Term{isomorphism}{同构}
\[c(m,n):\operatorname{loc}(n)\operatorname{loc}(m)\simeq\operatorname{loc}(nm),\tag{7.3.1.5}\]
且这些 $c(m,n)$ 满足与 (6.1.2.2) 类似的\Term{cocycle}{上循环}性质。
\end{SGAresult}
\begin{SGAresult}{例}{7.3.2}
（一个景的\Term{morphism}{态射}\Term{fibered-site}{纤维景}。）设 $C$ 为\Term{fibre-product}{纤维积}\Term{representable}{可表}的一个 $\mathcal U$-\Term{site}{景}，$\operatorname{Fl}(C)$ 为 $C$ 的\Term{morphism}{态射}\Term{category}{范畴}，$p:\operatorname{Fl}(C)\to C$ 将\Term{morphism}{态射}映到其终点。在\Term{fiber}{纤维}\Term{category}{范畴}上赋予\Term{induced-topology}{诱导拓扑}（第 II 讲 5.2）后，$(\operatorname{Fl}(C),p)$ 是 $C$ 上的一个 $\mathcal U$-\Term{fibered-site}{纤维景}。设 $\varepsilon_C:C\to C^\sim$ 为 $C$ 到其 $\mathcal U$-\Term{sheaf}{层}\Term{category}{范畴}的\Term{canonical}{典范}\Term{functor}{函子}。以 $\operatorname{Fl}(\varepsilon_C)$ 表示它到\Term{morphism}{态射}\Term{category}{范畴}的自然延伸，有\Term{category}{范畴}与\Term{functor}{函子}的\Term{commutative-diagram}{交换图}
\[\begin{tikzcd}
\operatorname{Fl}(C)\arrow[r,"\operatorname{Fl}(\varepsilon_C)"]\arrow[d,"p"']&\operatorname{Fl}(C^\sim)\arrow[d,"p"]\\
C\arrow[r,"\varepsilon_C"']&C^\sim.
\end{tikzcd}\tag{7.3.2.1}\]
由此得到 $C$ 上的\Term{fibered-category}{纤维范畴}间的一个\Term{cartesian-functor}{笛卡尔函子}
\[\begin{tikzcd}
\operatorname{Fl}(C)\arrow[rr]\arrow[dr,"p"']&&\operatorname{Fl}(C^\sim)_C\arrow[dl,"p_{/C}"]\\
&C&
\end{tikzcd}\tag{7.3.2.2}\]
其中 $p_{/C}:\operatorname{Fl}(C^\sim)_C\to C$ 由 $\operatorname{Fl}(C^\sim)\to C^\sim$ 经\Term{base-change}{基变换} $\varepsilon_C$ 得到。(7.3.2.2) 中的\Term{functor}{函子} $\operatorname{Fl}(C)\to\operatorname{Fl}(C^\sim)_C$ 是从\Term{fibered-site}{纤维景} $\operatorname{Fl}(C^\sim)_C$ 到 $\operatorname{Fl}(C)$ 的一个\Term{morphism}{态射}（7.2.2）。因此由 7.2.7 得到 $C$ 上的\Term{fibered-topos}{纤维拓扑斯}\Term{morphism}{态射}
\[\operatorname{can}:\operatorname{Fl}(C^\sim)_C\longrightarrow\operatorname{Fl}(C)^{\sim/C},\tag{7.3.2.3}\]
其中 $\operatorname{Fl}(C)^{\sim/C}$ 为与\Term{fibered-site}{纤维景} $\operatorname{Fl}(C)\to C$ 相伴的\Term{fibered-topos}{纤维拓扑斯}（7.2.6）。由相伴\Term{fibered-topos}{纤维拓扑斯}的定义及第 II 讲 5.5 可知，(7.3.2.3) 中的 $\operatorname{can}$ 是 $C$ 上 $\mathcal U$-\Term{fibered-topos}{纤维拓扑斯}的一个 $C$-\Term{equivalence-of-categories}{等价}，即存在一个 $C$ 上 $\mathcal U$-\Term{fibered-topos}{纤维拓扑斯}\Term{morphism}{态射}
\SourcePage{193}
$\Theta:\operatorname{Fl}(C)^{\sim/C}\to\operatorname{Fl}(C^\sim)_C$，使复合 $\Theta\circ\operatorname{can}$ 和 $\operatorname{can}\circ\Theta$ 都 $C$-\Term{isomorphism}{同构}于恒等。

设 $\phi:I\to C$ 为一个\Term{functor}{函子}。经\Term{base-change}{基变换}得到 $I$ 上的\Term{fibered-site}{纤维景}
\[p_I:\operatorname{Fl}(C)_I\longrightarrow I.\tag{7.3.2.4}\]
由于相伴\Term{fibered-topos}{纤维拓扑斯}的形成与\Term{base-change}{基变换}交换（7.2.6），有 $\mathcal U$-\Term{fibered-topos}{纤维拓扑斯}的 $I$-\Term{equivalence-of-categories}{等价}
\[\operatorname{can}_I:\operatorname{Fl}(C^\sim)_I\longrightarrow\operatorname{Fl}(C)_I^{\sim/I},\tag{7.3.2.5}\]
其中 $\operatorname{Fl}(C^\sim)_I\to I$ 由\Term{base-change}{基变换} $\varepsilon_C\circ\phi:I\to C^\sim$ 得到。

最后设 $\phi_1,\phi_2:I\to C$ 为两个\Term{functor}{函子}，$m:\phi_1\to\phi_2$ 为\Term{functor}{函子}\Term{morphism}{态射}。以 $\operatorname{Fl}(C)_1$、$\operatorname{Fl}(C)_2$ 分别表示 $p:\operatorname{Fl}(C)\to C$ 经\Term{base-change}{基变换} $\phi_1$、$\phi_2$ 得到的 $I$ 上\Term{fibered-site}{纤维景}。如同 7.3.1，\Term{constructible}{可构造} $I$ 上的\Term{fibered-site}{纤维景}\Term{morphism}{态射}
\[\operatorname{loc}(m):\operatorname{Fl}(C)_1\longrightarrow\operatorname{Fl}(C)_2.\tag{7.3.2.6}\]
设 $\operatorname{Fl}(C^\sim)_1\to I$、$\operatorname{Fl}(C^\sim)_2\to I$ 分别经\Term{base-change}{基变换} $\varepsilon_C\circ\phi_1$、$\varepsilon_C\circ\phi_2$ 得到。存在一个精确到\Term{isomorphism}{同构}交换的 $I$ 上\Term{morphism}{态射}图
\[\begin{tikzcd}
\operatorname{Fl}(C^\sim)_1\arrow[r,"\operatorname{loc}(\varepsilon_C*m)"]\arrow[d]&\operatorname{Fl}(C^\sim)_2\arrow[d]\\
\operatorname{Fl}(C)_1\arrow[r,"\operatorname{loc}(m)"']&\operatorname{Fl}(C)_2,
\end{tikzcd}\tag{7.3.2.7}\]
其中竖直\Term{morphism}{态射}由 (7.3.2.2) 中的\Term{morphism}{态射}经\Term{base-change}{基变换}得到。
\end{SGAresult}
\begin{SGAresult}{练习}{7.3.3}
（小\Term{fibered-site}{纤维景}与大\Term{fibered-site}{纤维景}。）本练习续第 IV 讲练习 4.10.6，采用其中的假设和记号。以 $M$ 表示由 $\mathcal M$ 中\Term{morphism}{态射}定义的 $\operatorname{Fl}(S)$ 的\Term{full-subcategory}{满子范畴}，以 $p:M\to S$ 表示将\Term{morphism}{态射}映到终点的\Term{functor}{函子}。

$6^\circ)$ 证明 $(M,p)$ 是 $S$ 上的一个\Term{fibered-site}{纤维景}；并证明当 $S$ 中\Term{fibre-product}{纤维积}\Term{representable}{可表}时，包含\Term{functor}{函子} $M\to\operatorname{Fl}(S)$ 是从\Term{fibered-site}{纤维景} $\operatorname{Fl}(S)\to S$ 到\Term{fibered-site}{纤维景} $M\to S$ 的一个\Term{cartesian-morphism}{笛卡尔态射}。

$7^\circ)$ 设 $M^{\sim/S}\to S$ 为与 $M\to S$ 相伴的\Term{fibered-topos}{纤维拓扑斯}。采用 7.3.2 的记号，证明存在 $S$ 上\Term{fibered-topos}{纤维拓扑斯}的一个\Term{canonical}{典范}\Term{morphism}{态射} $g:\operatorname{Fl}(S^\sim)_{/S}\to M^{\sim/S}$，在每个 $X\in\operatorname{ob}(S)$ 处的\Term{fiber}{纤维}上诱导\Term{morphism}{态射} $(\operatorname{Prol}_X,\operatorname{Res}_X)$（第 IV 讲 4.10.6，$4^\circ)$）。证明：尽管对于每个 $X\in\operatorname{ob}(S)$，\Term{morphism}{态射} $g_X:S_{/X}^\sim\to S(X)^\sim$ 有一个\Term{section}{截面}（同上），$g$ 一般却没有\Term{section}{截面}。
\end{SGAresult}
\SourcePage{194}
\textbf{7.4. \Term{fibered-topos}{纤维拓扑斯}或\Term{fibered-site}{纤维景}的\Term{total-topology}{全拓扑}。}
\begin{SGAresult}{定义}{7.4.1}
设 $p:C\to I$ 为一个 $\mathcal U$-\Term{fibered-site}{纤维景}，对于每个 $i\in\operatorname{ob}(I)$，以 $\alpha_{i!}:C_i\to C$ 表示\Term{fiber}{纤维}\Term{category}{范畴} $C_i$ 到 $C$ 的包含\Term{functor}{函子}。使所有 $\alpha_{i!}$（$i\in\operatorname{ob}(I)$）\Term{continuous-functor}{连续}的 $C$ 上最粗\Term{topology}{拓扑}（第 III 讲 3.6），称为 $C$ 上的\Term{total-topology}{全拓扑}。赋予\Term{total-topology}{全拓扑}的\Term{category}{范畴} $C$ 称为\Term{total-site}{全景}。
\end{SGAresult}
\begin{SGAresult}{命题}{7.4.2}
设 $p:C\to I$ 为一个 $\mathcal U$-\Term{fibered-site}{纤维景}，$T$ 为 $C$ 上的\Term{total-topology}{全拓扑}，$T_i$ 为每个 $i\in\operatorname{ob}(I)$ 处\Term{fiber}{纤维} $C_i$ 的\Term{topology}{拓扑}。

1) 设 $X$ 是 $C$ 中位于 $i$ 上方的\Term{object}{对象}。一个族 $(X_\beta\to X)_{\beta\in B}$ 对于 $T$ 为\Term{covering-family}{覆盖族}，当且仅当它可被 $C_i$ 中的一个 $T_i$-\Term{covering-family}{覆盖族} $(Y_\gamma\to X)_{\gamma\in\Gamma}$ \Term{refinement}{加细}。

2) \Term{topology}{拓扑} $T$ 在各\Term{category}{范畴} $C_i$ 上诱导\Term{topology}{拓扑} $T_i$。

3) 对每个 $i\in\operatorname{ob}(I)$，\Term{functor}{函子} $\alpha_{i!}:C_i\to C$ 都\Term{cocontinuous-functor}{余连续}。
\end{SGAresult}
\textbf{7.4.2.1.} 性质 3) 由 1) 及第 III 讲 2.1 得到。证明 1)。对 $C$ 的每个\Term{object}{对象} $X$，以 $J(X)$ 表示 1) 中所描述的 $C_{/X}$ 的\Term{sieve}{筛}的\Term{set}{集合}。可验证各 $J(X)$ 满足第 II 讲 1.1 的性质 T 1)、T 2)、T 3，因而定义 $C$ 上的一个\Term{topology}{拓扑} $T'$。$T'$ 是使各 $T_i$ 的\Term{covering-family}{覆盖族}仍为\Term{covering-family}{覆盖族}的最粗\Term{topology}{拓扑}。所以 $T'$ 不细于 $T$（第 III 讲 1.6）。为了证明 $T'=T$，只须证明：赋予 $C$ 以 $T'$ 时，每个 $\alpha_{i!}:C_i\to C$ 都\Term{continuous-functor}{连续}；换言之，只须证明 $T'$ 在 $C_i$ 上诱导的\Term{topology}{拓扑}就是 $T_i$，这样同时也证明了 2)。设 $T_i'$ 为 $T'$ 在 $C_i$ 上的\Term{induced-topology}{诱导拓扑}，$u:R\hookrightarrow X$ 为一个 $T_i'$-\Term{covering-sieve}{覆盖筛}。由第 III 讲 3.2，\Term{morphism}{态射} $\alpha_{i!}(u):\alpha_{i!}(R)\to\alpha_{i!}(X)$ 是\Term{bicovering}{双覆盖}的，特别是\Term{covering}{覆盖}的；依 $T'$ 的定义，这就是说 $R\hookrightarrow X$ 是 $T_i$-\Term{covering-sieve}{覆盖筛}。因此 $T_i'$ 不细于 $T_i$。为了证明 $T_i'$ 细于 $T_i$，依\Term{induced-topology}{诱导拓扑}的定义（第 III 讲 3.1），只须证明当 $C_i$、$C$ 分别赋予 $T_i$、$T'$ 时，$\alpha_{i!}:C_i\to C$ \Term{continuous-functor}{连续}。总之，由第 III 讲 1.2、1.5，证明性质 1)、2) 只须证明下述引理。

\begin{SGAresult}{引理}{7.4.2.2}
采用 7.4.2 和 7.4.2.1 的假设与记号。设 $X$ 为 $C$ 中位于 $i$ 上方的\Term{object}{对象}，$R\xrightarrow uX$ 为 $\widehat C_i$ 的一个\Term{morphism}{态射}。下列性质\Term{equivalence}{等价}：

i) $R\to X$ 对 $T_i$ 是\Term{bicovering}{双覆盖}的。

ii) $\alpha_{i!}(u):\alpha_{i!}(R)\to\alpha_{i!}(X)$ 对 $T'$ 是\Term{bicovering}{双覆盖}的。
\end{SGAresult}
\begin{proof}
i) $\Rightarrow$ ii)：若 $R\to X$ 对 $T_i$ 是\Term{bicovering}{双覆盖}的，显然 $\alpha_{i!}(u)$ 对 $T'$ 为\Term{covering}{覆盖}。还须证明它是\Term{bicovering}{双覆盖}的（第 II 讲 5.2）。设 $Y$ 为 $C$ 中位于 $I$ 的\Term{object}{对象} $j$ 上方的\Term{object}{对象}，$m,n:Y\rightrightarrows\alpha_{i!}(R)$ 为 $\widehat C$ 的两个\Term{morphism}{态射}，满足
\SourcePage{195}
$\alpha_{i!}(u)m=\alpha_{i!}(u)n$，并令 $f=p(\alpha_{i!}(u)m)$。由于\Term{functor}{函子} $\alpha_{i!}$ 与\Term{inductive-limit}{归纳极限}交换（第 I 讲 5.4.3），有（第 I 讲 3.4）
\[\operatorname{Hom}_{\widehat C}(Y,\alpha_{i!}(R))\simeq\varinjlim_{C_i/R}\operatorname{Hom}_C(Y,\alpha_{i!}(\cdot)).\]
对每个 $Z\in\operatorname{ob}C_i$，依 6.1.1 的记号有
\[\operatorname{Hom}_C(Y,Z)\simeq\coprod_{w\in\operatorname{Hom}(j,i)}\operatorname{Hom}_w(Y,Z),\]
从而
\[\varinjlim_{C_i/R}\operatorname{Hom}(Y,\alpha_{i!}(\cdot))\simeq\coprod_{w\in\operatorname{Hom}(j,i)}\varinjlim_{C_i/R}\operatorname{Hom}_w(Y,\alpha_{i!}(\cdot)).\]
但对每个 $w\in\operatorname{Hom}(j,i)$，有\Term{inverse-image-functor}{逆像函子} $w^*:C_i\to C_j$，且 $\operatorname{Hom}_w(Y,\alpha_{i!}(\cdot))\simeq\operatorname{Hom}_{C_j}(Y,w^*(\cdot))$。仍以 $w^*:\widehat C_i\to\widehat C_j$ 表示其到\Term{presheaf}{预层}的延伸（依第 I 讲 5.4，本应记为 $(w^*)_!$），注意到该\Term{functor}{函子}与\Term{inductive-limit}{归纳极限}交换（第 I 讲 5.4.3），便得\Term{isomorphism}{同构}
\[\operatorname{Hom}_{\widehat C}(Y,\alpha_{i!}(R))=\coprod_{w\in\operatorname{Hom}(j,i)}\operatorname{Hom}_{\widehat C_j}(Y,w^*(R)).\]
在此\Term{isomorphism}{同构}下，$m,n:Y\rightrightarrows\alpha_{i!}(R)$ 对应于两个\Term{morphism}{态射} $m',n':Y\rightrightarrows f^*R$，且 $f^*(u)m'=f^*(u)n'$。由于\Term{inverse-image-functor}{逆像函子}是一个\Term{topos-morphism}{拓扑斯态射}，特别地它\Term{continuous-functor}{连续}，所以 $f^*(u)$ 是\Term{bicovering}{双覆盖}的（第 III 讲 3.2）。故存在 $C_j$ 的一个\Term{covering-family}{覆盖族} $(Y_\delta\xrightarrow{v_\delta}Y)_\delta$，使每个 $\delta$ 都有 $m'v_\delta=n'v_\delta$。于是 $mv_\delta=nv_\delta$，故 $\alpha_{i!}(u)$ 对 $T'$ 是\Term{bicovering}{双覆盖}的。

ii) $\Rightarrow$ i)：假设 $\alpha_{i!}(R)\xrightarrow{\alpha_{i!}(u)}\alpha_{i!}(X)$ 是\Term{bicovering}{双覆盖}的。\Term{functor}{函子} $\alpha_{i!}:C_i\to C$ \Term{cocontinuous-functor}{余连续}（第 III 讲 2.1）。因而 $\alpha_i^*\alpha_{i!}(u)$ 是\Term{bicovering}{双覆盖}的（第 III 讲 2.3.2 及第 II 讲 5.3 ii)）。$\alpha_{i!}:C_i\to C$ 具有第 I 讲 5.14 的性质 (PPF)，故在 $\widehat C_i$ 中有\Term{cartesian-diagram}{笛卡尔图}（第 I 讲 5.14，3)）
\[\begin{tikzcd}
R\arrow[r]\arrow[d,"u"']&\alpha_i^*\alpha_{i!}(R)\arrow[d,"\alpha_i^*\alpha_{i!}(u)"]\\
X\arrow[r]&\alpha_i^*\alpha_{i!}(X).
\end{tikzcd}\]
所以 $u$ 是\Term{bicovering}{双覆盖}的（第 II 讲 5.2）。
\end{proof}
\begin{SGAresult}{注记}{7.4.3}
1) 命题 7.4.2 可在比\Term{fibered-site}{纤维景}稍为一般的框架中陈述和证明，但看来没有什么意义：将\Term{fibered-site}{纤维景} $p:C\to I$ 换成一个\Term{fibered-category}{纤维范畴} $p:C\to I$，其\Term{fiber}{纤维}配有\Term{topology}{拓扑}，\Term{inverse-image-functor}{逆像函子}对于这些\Term{topology}{拓扑}\Term{continuous-functor}{连续}即可
\SourcePage{196}
（而在\Term{fibered-site}{纤维景}的情形，还要求这些\Term{inverse-image-functor}{逆像函子}是\Term{site-morphism}{景态射}，第 IV 讲 4.9）。

2) 可以证明，\Term{fibered-site}{纤维景}的\Term{total-topology}{全拓扑}也是使各 $\alpha_{i!}:C_i\to C$ \Term{cocontinuous-functor}{余连续}的最细\Term{topology}{拓扑}（第 III 讲 2）。

3) 设 $p:C\to I$ 是一个 $\mathcal U$-\Term{fibered-site}{纤维景}，底\Term{category-equivalence}{范畴等价}于某个 $\mathcal U$-\Term{small-category}{小范畴}。那么\Term{total-site}{全景} $C$ 是一个 $\mathcal U$-\Term{site}{景}（第 II 讲 3.0.2）。事实上，若对 $I$ 的每个\Term{object}{对象} $i$，$(X_\beta)_{\beta\in B_i}$ 是 $C_i$ 的一个小\Term{topologically-generating-family}{拓扑生成族}，则 $(X_\beta)_{\beta\in\coprod_iB_i}$ 是 $C$ 的一个 $\mathcal U$-小\Term{topologically-generating-family}{拓扑生成族}。\Term{total-site}{全景}上的\Term{sheaf}{层}\Term{topos}{拓扑斯}称为\Term{total-topos}{全拓扑斯}，记为 $\operatorname{Top}(C)$。

4) 设 $p:C\to I$ 为\Term{small-category}{小范畴}上的一个 $\mathcal U$-\Term{fibered-site}{纤维景}。由于 $\alpha_{i!}:C_i\to C$ \Term{cocontinuous-functor}{余连续}，它产生一个从 $C_i^\sim$ 到 $\operatorname{Top}(C)$ 的\Term{topos-morphism}{拓扑斯态射} $\alpha_i=(\alpha_i^*,\alpha_{i*})$（第 IV 讲 4.7）。又由于 $\alpha_{i!}$ \Term{continuous-functor}{连续}，\Term{functor}{函子} $\alpha_i^*:\operatorname{Top}(C)\to C_i^\sim$ 有一个\Term{left-adjoint}{左伴随函子}，仍记为 $\alpha_{i!}:C_i^\sim\to\operatorname{Top}(C)$，它延伸 $\alpha_{i!}:C_i\to C$（第 III 讲 1.2 iv)）。
\end{SGAresult}
\begin{SGAresult}{命题}{7.4.4}
设 $p:C\to I$ 为\Term{small-category}{小范畴} $I$ 上的一个 $\mathcal U$-\Term{fibered-site}{纤维景}。

1) $C$ 上的一个\Term{presheaf}{预层} $F$ 对\Term{total-topology}{全拓扑}是\Term{sheaf}{层}，当且仅当对 $I$ 的每个\Term{object}{对象} $i$，\Term{presheaf}{预层} $F\circ\alpha_{i!}=F|C_i$ 是 $C_i$ 上的\Term{sheaf}{层}。

2) 从 $C$ 到\Term{site}{景} $D$ 的一个\Term{functor}{函子} $m$ 是\Term{total-site}{全景} $C$ 到\Term{site}{景} $D$ 的\Term{continuous-functor}{连续函子}，当且仅当对 $I$ 的每个\Term{object}{对象} $i$，$m|C_i=m\circ\alpha_{i!}:C_i\to D$ \Term{continuous-functor}{连续}。
\end{SGAresult}
\begin{proof}
若 $F$ 对\Term{total-topology}{全拓扑}是\Term{sheaf}{层}，则 $F\circ\alpha_{i!}$ 是 $C_i$ 上的\Term{sheaf}{层}，因为 $\alpha_{i!}:C_i\to C$ \Term{continuous-functor}{连续}。反之，假设每个 $F\circ\alpha_{i!}$ 都是 $C_i$ 上的\Term{sheaf}{层}。设 $X$ 为 $C$ 中位于 $i$ 上方的\Term{object}{对象}。对 $C_{i/X}$ 中每个\Term{covering-sieve}{覆盖筛} $R$，有 $F(X)\simeq F\circ\alpha_{i!}(R)$，因而 $F(X)\simeq F(\alpha_{i!}(R))$。设 $R'\hookrightarrow X$ 是\Term{canonical}{典范}\Term{presheaf}{预层}\Term{morphism}{态射} $\alpha_{i!}(R)\to X$ 的像\Term{sieve}{筛}。$\alpha_{i!}(R)\to R'$ 是\Term{presheaf}{预层}的一个\Term{epimorphism}{满态射}，所以 $F(R')\to F(\alpha_{i!}(R))$ 是\Term{injectivity}{单射}。它同时也是\Term{surjectivity}{满射}，因而是\Term{bijection}{双射}，故 $F(X)\to F(R')$ 是\Term{bijection}{双射}。由于这些\Term{sieve}{筛}\Term{covering}{覆盖} $X$（7.4.2，1)），\Term{presheaf}{预层} $F$ 是\Term{sheaf}{层}：回顾\Term{associated-sheaf}{相伴层}的构造（第 II 讲 3），便知 $F$ \Term{isomorphism}{同构}于其\Term{associated-sheaf}{相伴层}（第 II 讲 3.3）。第二个断言立即由第一个断言和\Term{continuity}{连续性}的定义（第 III 讲 1.1）得到。
\end{proof}
\textbf{7.4.5.} 设 $p:C\to I$ 是一个 $\mathcal U$-\Term{fibered-site}{纤维景}，$I$ \Term{equivalence}{等价}于\Term{small-category}{小范畴}，$f:i\to j$ 为\Term{index-category}{指标范畴} $I$ 的一支\Term{arrow}{箭头}。仍以 $f:C_i^\sim\to C_j^\sim$ 表示\Term{fibered-site}{纤维景}结构所确定的\Term{topos-morphism}{拓扑斯态射}，采用 7.4.3，4) 的记号，有一个 $\mathcal U$-\Term{topos}{拓扑斯}\Term{morphism}{态射}图
\SourcePage{197}
\[\begin{tikzcd}
C_i^\sim\arrow[r,"\alpha_i"]\arrow[d,"f"']&\operatorname{Top}(C)\\
C_j\arrow[ur,"\alpha_j"']&
\end{tikzcd}\tag{7.4.5.1}\]
其中 $\operatorname{Top}(C)$ 为\Term{total-topos}{全拓扑斯}（7.4.3，3)）。(7.4.5.1) 一般不交换。下面定义一个\Term{canonical}{典范}的\Term{topos-morphism}{拓扑斯态射}间\Term{morphism}{态射}
\[\rho_f:\alpha_i\longrightarrow\alpha_j\circ f.\tag{7.4.5.2}\]
只须在\Term{inverse-image}{逆像}上定义这样的\Term{morphism}{态射}，即定义一个\Term{functor}{函子}\Term{morphism}{态射}
\[\beta_f^*:f^*\circ\alpha_j^*\longrightarrow\alpha_i^*,\tag{7.4.5.3}\]
或者利用 $f^*$ 与 $f_*$ 的\Term{adjunction}{伴随}，只须定义
\[\gamma_f:\alpha_j^*\longrightarrow f_*\alpha_i^*.\tag{7.4.5.4}\]
设 $Y$ 为 $C_j$ 的一个\Term{object}{对象}，$F$ 为\Term{total-site}{全景} $C$ 上的一个\Term{sheaf}{层}。依定义，$\alpha_j^*F(Y)=F(\alpha_{j!}(Y))$ 且 $f_*\alpha_i^*F(Y)=F(\alpha_{i!}f^*(Y))$，其中右端的 $f^*:C_j\to C_i$ 表示\Term{fibered-structure}{纤维结构}的\Term{inverse-image-functor}{逆像函子}。依这个\Term{functor}{函子}的定义，有\Term{canonical}{典范}\Term{cartesian-morphism}{笛卡尔态射}
\[m_Y:\alpha_{i!}f^*\longrightarrow\alpha_{j!}.\tag{7.4.5.5}\]
应用\Term{functor}{函子} $F$，得到关于 $Y$ \Term{functorial}{函子化}的映射
\[\gamma_f(F)(Y)=F(m_Y):\alpha_j^*F(Y)\longrightarrow f_*\alpha_i^*F(Y),\tag{7.4.5.6}\]
它定义 $C_j$ 上的一个\Term{sheaf}{层}\Term{morphism}{态射}
\[\gamma_f(F):\alpha_j^*F\longrightarrow f_*\alpha_i^*F,\tag{7.4.5.7}\]
从而得到 (7.4.5.4) 的\Term{functor}{函子}\Term{morphism}{态射} $\gamma_f$。

设 $f:i\to j$ 和 $g:j\to k$ 为 $I$ 的两个可复合\Term{morphism}{态射}。有\Term{canonical}{典范}\Term{isomorphism}{同构} $c'_{g,f}:g_*f_*\to(gf)_*$，并可立即验证相容公式
\[c'_{g,f}\circ g_*(\gamma_f)\circ\gamma_g=\gamma_{gf}.\tag{7.4.5.8}\]
\SourcePage{198}
\textbf{7.4.6.} 现在引入与 $p:C\to I$ 相伴的 $\mathcal U$-\Term{fibered-topos}{纤维拓扑斯} $p^{\sim/I}:C^{\sim/I}\to I$（7.2.6），以及相应的 $I^\circ$ 上\Term{fibered-category}{纤维范畴} $(p^{\sim/I})':(C^{\sim/I})'\to I^\circ$（7.1.3.2）。以上讨论将 $\operatorname{Top}(C)$ 的每个\Term{object}{对象} $F$ 映到 $C_i^\sim$ 的\Term{object}{对象}族 $F_i=(\alpha_i^*F)_{i\in\operatorname{ob}(I)}$ 及\Term{morphism}{态射}族
\[\gamma_fF:F_j\longrightarrow f_*F_i,\qquad f\in\operatorname{Fl}(I),\quad f:i\to j,\]
这些\Term{morphism}{态射}满足 (7.4.5.8) 的相容条件。因此与 $F$ 相联系的是从 $I^\circ$ 到 $(C^{\sim/I})'$ 的一个 $I^\circ$-\Term{functor}{函子}，即 $\mathscr{H}\!om_{I^\circ}(I^\circ,(C^{\sim/I})')$ 的一个\Term{object}{对象} $\Theta_C(F)$。此构造关于 $F$ 具有\Term{functorial}{函子性}，故最终得到\Term{functor}{函子}
\[\Theta_C:\operatorname{Top}(C)\longrightarrow\mathscr{H}\!om_{I^\circ}(I^\circ,(C^{\sim/I})').\tag{7.4.6.1}\]
\begin{SGAresult}{命题}{7.4.7}
设 $C\to I$ 是一个 $\mathcal U$-\Term{fibered-site}{纤维景}，$I$ \Term{equivalence}{等价}于\Term{small-category}{小范畴}。\Term{functor}{函子} $\Theta_C$（7.4.6.1）是一个\Term{equivalence-of-categories}{范畴等价}。
\end{SGAresult}
\begin{proof}
只需描述 $\Theta_C$ 的一个\Term{quasi-inverse}{拟逆函子}。设 $\sigma\in\operatorname{Hom}_{I^\circ}(I^\circ,(C^\sim)')$，即 $i\mapsto\sigma(i)$ 是从 $I^\circ$ 到 $(C^{\sim/I})'$ 的一个 $I^\circ$-\Term{functor}{函子}。对于 $C$ 中位于 $p(X)=i\in\operatorname{ob}(I)$ 上方的每个\Term{object}{对象} $X$，有一个\Term{sheaf}{层} $\sigma(p(X))\in\operatorname{ob}C_i^\sim$，从而有\Term{set}{集合} $\sigma(p(X))(X)$。设 $m:X\to Y$ 为位于 $I$ 中 $f:i\to j$ 上方的\Term{morphism}{态射}。$m$ 唯一地\Term{resolution}{分解}为 $C_i$ 中的\Term{morphism}{态射} $m':X\to f^*Y$ 与\Term{canonical}{典范}\Term{cartesian-morphism}{笛卡尔态射} $f^*(Y)\to Y$ 的复合。同样，$\sigma(f)$ 唯一地\Term{resolution}{分解}为\Term{morphism}{态射} $\gamma_f(\sigma):\sigma(p(Y))\to f_*\sigma(p(X))$ 与\Term{canonical}{典范} $I^\circ$-\Term{cartesian-morphism}{笛卡尔态射} $f_*\sigma(p(X))\to\sigma(p(X))$ 的复合。因此复合下列映射便得到从 $\sigma(p(Y))(Y)$ 到 $\sigma(p(X))(X)$ 的映射：
\[\begin{aligned}
\sigma(p(Y))(Y)&\xrightarrow{\gamma_f(\sigma)(Y)}f_*(p(X))(Y)\\
&=\sigma(p(X))(f^*Y)\xrightarrow{\sigma(p(X))(m')}\sigma(p(X))(X).
\end{aligned}\]
这样便将 $C$ 的每个\Term{object}{对象} $X$ 映到\Term{set}{集合} $\sigma(p(X))(X)$，将每个\Term{morphism}{态射} $m:X\to Y$ 映到映射 $\sigma(p(X))(Y)\to\sigma(p(X))(X)$。可验证这确实定义 $C$ 上的一个\Term{presheaf}{预层}，且不依赖于构造中\Term{inverse-image-functor}{逆像函子}的选择。由 7.4.4，$X\mapsto\sigma(p(X))(X)$ 是\Term{total-site}{全景} $C$ 上的\Term{sheaf}{层}，而且显然关于 $\mathscr{H}\!om_{I^\circ}(I^\circ,(C^{\sim/I})')$ 的\Term{object}{对象} $\sigma$ 具有\Term{functorial}{函子性}。回顾 $\Theta_C$ 的定义，也显然可知刚构造的从 $\mathscr{H}\!om_{I^\circ}(I^\circ,(C^{\sim/I})')$ 到 $\operatorname{Top}(C)$ 的\Term{functor}{函子}是 $\Theta_C$ 的一个\Term{quasi-inverse}{拟逆函子}。
\end{proof}
\begin{SGAresult}{注记}{7.4.8}
特别地，由 7.4.7 可知 $\mathscr{H}\!om_{I^\circ}(I^\circ,(C^{\sim/I})')$ 是一个 $\mathcal U$-\Term{topos}{拓扑斯}。这也可直接看出：利用第 I 讲 9.21.10，立即可知第 IV 讲 1.1.2 的条件 a)、b)、c) 成立；小\Term{generating-category}{生成范畴}的存在性来自第 I 讲 9.25。
\end{SGAresult}
\SourcePage{199}
\textbf{7.4.9.} 设 $p:C\to I$、$q:D\to I$ 是\Term{small-category}{小范畴}上的两个 $\mathcal U$-\Term{fibered-site}{纤维景}，$m$ 是从\Term{fibered-site}{纤维景} $C$ 到 $D$ 的一个\Term{morphism}{态射}（7.2.2）。以 $m^{\sim/I}:C^{\sim/I}\to D^{\sim/I}$ 表示相伴 $\mathcal U$-\Term{fibered-topos}{纤维拓扑斯}间对应的\Term{morphism}{态射}；它确定到\Term{canonical}{典范}\Term{isomorphism}{同构}（7.2.7）。相应的\Term{direct-image-functor}{正像函子} $m_*^{\sim/I}:(C^{\sim/I})'\to(D^{\sim/I})'$（7.1.7）在 $I^\circ$ 上\Term{section}{截面}\Term{category}{范畴}上给出\Term{functor}{函子}
\[\begin{aligned}
\operatorname{Hom}_{I^\circ}(I^\circ,m_*^{\sim/I}):\mathscr{H}\!om_{I^\circ}(I^\circ,(C^{\sim/I})')
\longrightarrow\mathscr{H}\!om_{I^\circ}(I^\circ,(D^{\sim/I})').
\end{aligned}\tag{7.4.9.1}\]
另一方面，$m$ 是\Term{total-site}{全景}之间的\Term{continuous-functor}{连续函子}（7.4.4），所以通过复合得到\Term{functor}{函子} $m_*:\operatorname{Top}(C)\to\operatorname{Top}(F)$。总之，有\Term{functor}{函子}图
\[\begin{tikzcd}[column sep=large]
\operatorname{Top}(C)\arrow[r,"\Theta_C"]\arrow[d,"m_*"']&\mathscr{H}\!om_{I^\circ}(I^\circ,(C^{\sim/I})')\arrow[d,"{\operatorname{Hom}_{I^\circ}(I^\circ,m_*^{\sim/I})}"]\\
\operatorname{Top}(D)\arrow[r,"\Theta_D"']&\mathscr{H}\!om_{I^\circ}(I^\circ,(D^{\sim/I})').
\end{tikzcd}\tag{7.4.9.2}\]
\begin{SGAresult}{命题}{7.4.10}
图 (7.4.9.2) 精确到\Term{isomorphism}{同构}交换。若 $m$ 是\Term{cartesian-functor}{笛卡尔}的，则 $m$ 是从\Term{total-site}{全景} $C$ 到\Term{total-site}{全景} $D$ 的一个\Term{site-morphism}{景态射}。
\end{SGAresult}
\begin{proof}
图 (7.4.9.2) 精确到\Term{isomorphism}{同构}的交换性留给读者验证。考虑到 7.4.7，为证明 $m$ 是\Term{total-site}{全景}之间的\Term{morphism}{态射}，只须证明 $\operatorname{Hom}_{I^\circ}(I^\circ,m_*^{\sim/I})$ 的\Term{left-adjoint}{左伴随函子}\Term{left-exact}{左正合}。先描述\Term{morphism}{态射} $m^{\sim/I}$ 的\Term{inverse-image-functor}{逆像函子}（7.2.7），对每个\Term{object}{对象} $i$，以 $m_i^*:D_i^\sim\to C_i^\sim$ 表示 $(m^{\sim/I})^*$ 在 $i$ 处\Term{fiber}{纤维}上诱导的\Term{functor}{函子}。对每个 $f:i\to j$，在两边分别选取\Term{topos-morphism}{拓扑斯态射} $(f^*,f_*):C_i^\sim\to C_j^\sim$ 和 $(f^*,f_*):D_i^\sim\to D_j^\sim$。因为 $(m^{\sim/I})^*$ 是一个 $I^\circ$-\Term{cartesian-functor}{笛卡尔函子}，每个 $f:i\to j$ 给出\Term{canonical}{典范}\Term{isomorphism}{同构}
\[m_i^*f_*\longrightarrow f^*m_j^*,\tag{7.4.10.1}\]
而 $f^*$ 与 $f_*$ \Term{adjunction}{伴随}，故有\Term{canonical}{典范}\Term{morphism}{态射}（\Term{adjunction}{伴随}\Term{morphism}{态射}）
\[\phi:\operatorname{id}\longrightarrow f_*f^*,\qquad\psi:f^*f_*\longrightarrow\operatorname{id}.\tag{7.4.10.2}\]
考虑下列\Term{functor}{函子}\Term{morphism}{态射}序列：
\[m_j^*f^*\xrightarrow{(1)}f_*f^*m_j^*f_*\xrightarrow{(2)}f_*m_i^*f^*f_*\xrightarrow{(3)}f_*m_i^*.\tag{7.4.10.3}\]
\SourcePage{200}
其中 (1) 由 $\phi$ 得到，(2) 由\Term{isomorphism}{同构} (7.4.10.1) 得到，(3) 由 $\psi$ 得到。将 (7.4.10.3) 中各\Term{morphism}{态射}复合，得到\Term{functor}{函子}\Term{morphism}{态射}
\[\operatorname{can}_f:m_j^*f_*\longrightarrow f_*m_i^*.\tag{7.4.10.4}\]
现设 $\tau$ 为 $\mathscr{H}\!om_{I^\circ}(I^\circ,(D^{\sim/I})')$ 的\Term{object}{对象}。对每个 $i$ 有\Term{object}{对象} $\tau(i)\in D_i^\sim$，对每个 $f:i\to j$ 有\Term{morphism}{态射}
\[\gamma_f(\tau):\tau(j)\longrightarrow f_*\tau(i),\tag{7.4.10.5}\]
且 $\gamma_f(\tau)$ 族满足与 (7.4.5.8) 类似的相容性。置 $\sigma(i)=m_i^*(\tau(i))$，并以
\[\gamma_f(\sigma):\sigma(j)\longrightarrow f_*\sigma(i)\tag{7.4.10.6}\]
表示复合 $m_j^*\tau(j)\xrightarrow{m_j^*\gamma_f(\tau)}m_j^*f_*\tau(i)\xrightarrow{\operatorname{can}_f(\tau(i))}f_*m_i^*\tau(i)$。可验证这些 $\gamma_f(\sigma)$ 满足 (7.4.5.8) 的相容性，故定义 $\mathscr{H}\!om_{I^\circ}(I^\circ,(C^{\sim/I})')$ 的一个\Term{object}{对象}。作一番“\Term{adjoint-functors}{伴随函子}追踪”便知，刚构造的从 $\mathscr{H}\!om_{I^\circ}(I^\circ,(D^{\sim/I})')$ 到 $\mathscr{H}\!om_{I^\circ}(I^\circ,(C^{\sim/I})')$ 的\Term{functor}{函子}\Term{left-adjoint}{左伴随}于 $\operatorname{Hom}_{I^\circ}(I^\circ,m_*^{\sim/I})$。还须证明此\Term{left-adjoint}{左伴随函子}与\Term{finite}{有限}\Term{projective-limit}{投射极限}交换；这立即来自\Term{section}{截面}\Term{category}{范畴}中的\Term{projective-limit}{投射极限}可逐\Term{fiber}{纤维}计算，以及构造该\Term{left-adjoint}{左伴随函子}所用的所有\Term{functor}{函子}都与\Term{finite}{有限}\Term{projective-limit}{投射极限}交换。
\end{proof}
\begin{SGAresult}{推论}{7.4.11}
设 $p:C\to I$ 是\Term{small-category}{小范畴}上的一个 $\mathcal U$-\Term{fibered-site}{纤维景}，$p^{\sim/I}:C^{\sim/I}\to I$ 是其相伴 $\mathcal U$-\Term{fibered-topos}{纤维拓扑斯}（7.2.6），$\varepsilon_I:C\to C^{\sim/I}$ 是从 $\mathcal U$-\Term{fibered-site}{纤维景} $C^{\sim/I}$ 到 $C$ 的\Term{canonical}{典范}\Term{morphism}{态射}。$\varepsilon_I$ 在相应\Term{total-topos}{全拓扑斯}之间诱导一个\Term{equivalence-of-categories}{等价} $\operatorname{Top}(\varepsilon_I):\operatorname{Top}(C)\to\operatorname{Top}(C^{\sim/I})$。
\end{SGAresult}
\begin{proof}
由 (7.4.9.2) 的交换性，以及 $\varepsilon_I$ 在相伴 $\mathcal U$-\Term{fibered-topos}{纤维拓扑斯}之间给出\Term{equivalence-of-categories}{等价}这一事实得到。
\end{proof}
\begin{SGAresult}{注记}{7.4.11.1}
由 7.4.11，在关于\Term{total-topos}{全拓扑斯}的问题中，可用相应的\Term{fibered-topos}{纤维拓扑斯}代替\Term{fibered-site}{纤维景}。
\end{SGAresult}
\begin{SGAresult}{引理}{7.4.12}
设 $p:C\to I$ 为\Term{small-category}{小范畴} $I$ 上的一个\Term{fibered-site}{纤维景}，$i$ 为 $I$ 的一个\Term{final-object}{终对象}（第 I 讲 10）。\Term{functor}{函子} $\alpha_{i!}:C_i^\sim\to\operatorname{Top}(C)$（7.4.3，4)）\Term{exact-functor}{正合}且\Term{fully-faithful}{完全忠实}。\Term{functor}{函子}对 $\beta_i=(\alpha_{i!},\alpha_i^*):\operatorname{Top}(C)\to C_i^\sim$ 和 $\alpha_i=(\alpha_i^*,\alpha_{i*}):C_i^\sim\to\operatorname{Top}(C)$ 是\Term{topos-morphism}{拓扑斯态射}。复合\Term{morphism}{态射} $\beta_i\alpha_i:C_i^\sim\to C_i^\sim$ \Term{isomorphism}{同构}于恒等\Term{morphism}{态射}。
\end{SGAresult}
\begin{proof}
对 $I$ 的每个\Term{object}{对象} $j$，以 $f_j:j\to i$ 表示唯一的\Term{morphism}{态射}。利用\Term{equivalence-of-categories}{等价} $\Theta_C:\operatorname{Top}(C)\xrightarrow\sim\mathscr{H}\!om_{I^\circ}(I^\circ,(C^{\sim/I})')$（7.4.7）可知：对 $C_i$ 的每个\Term{object}{对象} $X$，$\alpha_{i!}(X)$ 是\Term{section}{截面} $j\mapsto f_j^*(X)$；对每个\Term{section}{截面} $j\mapsto\sigma(j)\in\operatorname{Hom}_{I^\circ}(I^\circ,(C^{\sim/I})')$，$\alpha_i^*(j\mapsto\sigma(j))$ 是\Term{object}{对象} $\sigma(i)$。故 $\alpha_i^*\alpha_{i!}$ \Term{isomorphism}{同构}于恒等。
\SourcePage{201}
因此 $\alpha_{i!}$ \Term{fully-faithful}{完全忠实}。由于各 $f_j^*$ \Term{left-exact}{左正合}，$\alpha_{i!}$ \Term{left-exact}{左正合}。所以 $\beta_i$ 是\Term{topos-morphism}{拓扑斯态射}；又因 $\alpha_i^*\alpha_{i!}$ \Term{isomorphism}{同构}于恒等，$\beta_i\alpha_i$ \Term{isomorphism}{同构}于恒等\Term{morphism}{态射}。
\end{proof}
\begin{SGAresult}{定理}{7.4.13.1}
设 $p:C\to I$、$q:D\to I$ 是\Term{small-category}{小范畴} $I$ 上的两个 $\mathcal U$-\Term{fibered-site}{纤维景}，$m:D\to C$ 是一个 $I$-\Term{functor}{函子}。为了使 $m$ 是从\Term{fibered-site}{纤维景} $C$ 到 $D$ 的一个\Term{morphism}{态射}（7.2.2），只须 $m$ 是从\Term{total-site}{全景} $C$ 到\Term{total-site}{全景} $D$ 的一个\Term{site-morphism}{景态射}。
\end{SGAresult}
\textbf{7.4.13.2.} 必须证明：若 $I$-\Term{functor}{函子} $m:D\to C$ 是\Term{total-site}{全景} $C\to D$ 的\Term{site-morphism}{景态射}，则它是\Term{fibered-site}{纤维景} $C\to D$ 的\Term{morphism}{态射}，即每个\Term{object}{对象} $i$ 处的\Term{fiber}{纤维}上诱导的 $m_i:D_i\to C_i$ 都是\Term{site-morphism}{景态射} $C_i\to D_i$。此断言的证明占据 7.4.13.3 至 7.4.13.7。

\textbf{7.4.13.3.} 证明每个 $m_i:D_i\to C_i$ 都\Term{continuous-functor}{连续}。有\Term{commutative-diagram}{交换图}
\[\begin{tikzcd}
D_i\arrow[r,"m_i"]\arrow[d,"\alpha_{i!}"']&C_i\arrow[d,"\alpha_{i!}"]\\
D\arrow[r,"m"']&C.
\end{tikzcd}\]
只须证明：对于 $C_{i/X}$ 的每个\Term{covering-sieve}{覆盖筛} $R\hookrightarrow X$，\Term{morphism}{态射} $m_{i!}(R)\to m_i(X)$ 是\Term{bicovering}{双覆盖}的（第 III 讲 1.2 和 1.5）。由上图的交换性及 $\alpha_{i!}$、$m$ 的\Term{continuity}{连续性}，\Term{morphism}{态射} $\alpha_{i!}m_{i!}(R)\to\alpha_{i!}m_i(X)$ 是\Term{bicovering}{双覆盖}的。故所求断言由 7.4.2.2 得到。

\textbf{7.4.13.4. 归结为 $C,D$ 是 $\mathcal U$-\Term{fibered-topos}{纤维拓扑斯}的情形。} 因各 $m_i$ \Term{continuous-functor}{连续}，它们可自然延伸到\Term{sheaf}{层}\Term{category}{范畴}，记为 $m_i^*:D_i^\sim\to C_i^\sim$（第 III 讲 1.2 iv)）。设 $D^{\sim/I}$、$C^{\sim/I}$ 分别为与 $D$、$C$ 相伴的 $\mathcal U$-\Term{fibered-topos}{纤维拓扑斯}。对 $I$ 的每个\Term{morphism}{态射} $f$，以 $f^*$ 表示四个\Term{fibered-category}{纤维范畴} $C,D,C^{\sim/I},D^{\sim/I}$ 的\Term{inverse-image-functor}{逆像函子}。因 $m$ 是 $I$-\Term{functor}{函子}，对每个 $f:i\to j$ 有\Term{canonical}{典范}\Term{morphism}{态射}
\[f^*m_j\longleftarrow m_if^*.\]
由于 $f^*,m_i,m_j$ 都\Term{continuous-functor}{连续}，转到\Term{sheaf}{层}\Term{category}{范畴}便得\Term{canonical}{典范}\Term{isomorphism}{同构}
\[f^*m_j^*\longleftarrow m_i^*f^*.\]
这些\Term{isomorphism}{同构}满足相容性，可用来构造一个\Term{cartesian-functor}{笛卡尔函子} $m^*:D^{\sim/I}\to C^{\sim/I}$，在各\Term{fiber}{纤维}\Term{topos}{拓扑斯}上诱导 $m_i^*$。
\SourcePage{202}
此外，图
\[\begin{tikzcd}
D\arrow[r,"m"]\arrow[d,"\varepsilon_I"']&C\arrow[d,"\varepsilon_I"]\\
D^{\sim/I}\arrow[r,"m^*"']&C^{\sim/I}
\end{tikzcd}\]
精确到\Term{isomorphism}{同构}交换。由于各 $\varepsilon_I$ 在\Term{total-topos}{全拓扑斯}上诱导\Term{equivalence-of-categories}{等价}（7.4.11），且 $m$ 是\Term{total-site}{全景}之间的\Term{morphism}{态射}，$m^*:D^{\sim/I}\to C^{\sim/I}$ 也是\Term{total-site}{全景}之间的\Term{morphism}{态射}。故只须在 $D,C$ 为 $\mathcal U$-\Term{fibered-topos}{纤维拓扑斯}时证明 7.4.13.2；以下作此假设。

\textbf{7.4.13.5.} 对每个\Term{object}{对象} $i$，\Term{functor}{函子} $m_i:D_i\to C_i$ 将 $D_i$ 的\Term{final-object}{终对象}映到 $C_i$ 的\Term{final-object}{终对象}。以 $\overline m:D^\sim\to C^\sim$ 表示 $m$ 在\Term{total-topos}{全拓扑斯}上的自然延伸（第 III 讲 1.2）。利用\Term{equivalence-of-categories}{等价} $\Theta_D:D^\sim\to\mathscr{H}\!om_{I^\circ}(I^\circ,D')$（7.4.7），立即可验证，$\overline m$ 将每个\Term{section}{截面} $\sigma\in\operatorname{Hom}_{I^\circ}(I^\circ,D')$ 映到\Term{section}{截面} $(i\mapsto m_i\sigma(i))\in\operatorname{Hom}_{I^\circ}(I^\circ,C')$。因 $m$ 是\Term{site-morphism}{景态射}，$\overline m$ \Term{left-exact}{左正合}，特别地将 $\mathscr{H}\!om_{I^\circ}(I^\circ,D)$ 的\Term{final-object}{终对象}映到 $\mathscr{H}\!om_{I^\circ}(I^\circ,C')$ 的\Term{final-object}{终对象}。对每个 $i$，设 $e_i$ 为 $D_i$ 的一个\Term{final-object}{终对象}。\Term{section}{截面} $i\mapsto e_i$ 是 $\mathscr{H}\!om_{I^\circ}(I^\circ,D')$ 的一个\Term{final-object}{终对象}，所以 $i\mapsto m_i(e_i)$ 是 $\mathscr{H}\!om_{I^\circ}(I^\circ,C')$ 的一个\Term{final-object}{终对象}，故每个 $m_i(e_i)$ 都是 $C_i$ 的\Term{final-object}{终对象}。

\textbf{7.4.13.6. 归结为 $i$ 是 $I$ 的\Term{final-object}{终对象}的情形。} 设 $i$ 为 $I$ 的\Term{object}{对象}。\Term{functor}{函子} $\alpha_{i!}:D_i\to D$ \Term{resolution}{分解}为
\[D_i\xrightarrow{\overline\alpha_{i!}}D_{/e_i}\xrightarrow{\operatorname{loc}(e_i)}D,\]
其中 $\overline\alpha_{i!}$ 是显然的\Term{functor}{函子}，$\operatorname{loc}(e_i)$ 是\Term{forgetful-functor}{遗忘函子}。以 $m_{/i}:D_{/e_i}\to C_{/m(e_i)}$ 表示将 $D_{/e_i}$ 的\Term{object}{对象} $u:X\to e_i$ 映到 $C_{/m(e_i)}$ 的\Term{object}{对象} $m(u):m(X)\to m(e_i)$ 的\Term{functor}{函子}。有\Term{commutative-diagram}{交换图}
\[\begin{tikzcd}
D_i\arrow[r,"m_i"]\arrow[d,"\overline\alpha_{i!}"']&C_i\arrow[d,"\overline\alpha_{i!}"]\\
D_{/e_i}\arrow[r,"m_{/i}"']&C_{/m(e_i)}.
\end{tikzcd}\]
因 $e_i,m(e_i)$ 分别为 $D_i,C_i$ 的\Term{final-object}{终对象}（7.4.13.5），$D_{/e_i}$、$C_{/m(e_i)}$ 是 $I_{/i}$ 上的 $\mathcal U$-\Term{fibered-topos}{纤维拓扑斯}，而 $m_{/i}$ 是一个 $I_{/i}$-\Term{functor}{函子}。由于\Term{functor}{函子}为\Term{site-morphism}{景态射}这一性质可以\Term{localization}{局部化}（第 IV 讲 4.9），$m_{/i}$ 是从\Term{total-site}{全景} $C_{/m(e_i)}$ 到\Term{total-site}{全景}
\SourcePage{203}
$D_{/e_i}$ 的\Term{morphism}{态射}。故只须在 $i$ 是 $I$ 的\Term{final-object}{终对象}时证明 $m_i$ 是\Term{topos-morphism}{拓扑斯态射}；以下作此假设。

\textbf{7.4.13.7. 证明的结束。} 分别以 $\alpha_{i!}:D_i\to D^\sim$、$m^\sim:D^\sim\to C^\sim$、$\alpha_{i!}:C_i\to C^\sim$ 表示 $\alpha_{i!}$、$m$ 在\Term{sheaf}{层}\Term{category}{范畴}上的自然延伸。存在精确到\Term{isomorphism}{同构}交换的图
\[\begin{tikzcd}
D_i\arrow[r,"m_i"]\arrow[d,"\alpha_{i!}"']&C_i\arrow[d,"\alpha_{i!}"]\\
D^\sim\arrow[r,"m^\sim"']&C.
\end{tikzcd}\]
由于 $i$ 是\Term{final-object}{终对象}，$\alpha_i^*\alpha_{i!}:C_i\to C_i$ \Term{isomorphism}{同构}于恒等（7.4.12），所以 $m_i$ \Term{isomorphism}{同构}于 $\alpha_i^*m^\sim\alpha_{i!}$。又 $\alpha_i^*$、$m^\sim$、$\alpha_i^*$ 都\Term{left-exact}{左正合}（7.4.12）。所以 $m_i$ \Term{left-exact}{左正合}，因而是从\Term{site}{景} $C_i$ 到\Term{site}{景} $D_i$ 的一个\Term{morphism}{态射}。

\begin{SGAresult}{练习}{7.4.14}
设 $X=(X_i)_{i\in I}$ 为 $I$ 上的一个\Term{fibered-topos}{纤维拓扑斯}，$T$ 为一个\Term{topos}{拓扑斯}。证明各\Term{category}{范畴} $\mathscr{H}\!omtop(X_i,T)$ 是某个\Term{fibered-category}{纤维范畴}的\Term{fiber}{纤维}，将该\Term{category}{范畴}记为 $\mathscr{H}\!omtop_I(X,T\times I)$。证明它在 $I$ 上的\Term{section}{截面}\Term{category}{范畴}\Term{canonical}{典范}\Term{equivalence}{等价}于 $\mathscr{H}\!omtop_I(\operatorname{Top}(X),T)$，其中 $\operatorname{Top}(X)$ 是 $X$ 的\Term{total-topos}{全拓扑斯}。
\end{SGAresult}
\begin{SGAresult}{练习}{7.4.15}
设 $X=(X_i)_{i\in I}$ 为 $I$ 上的一个\Term{fibered-topos}{纤维拓扑斯}。定义从 $X$ 到\Term{fiber}{纤维}为 $(\mathrm{Ens})$ 的常值\Term{fibered-topos}{纤维拓扑斯}的一个\Term{cartesian-morphism}{笛卡尔态射} $m$。证明\Term{fiber}{纤维}为 $(\mathrm{Ens})$ 的常值\Term{fibered-topos}{纤维拓扑斯}的\Term{total-topos}{全拓扑斯}\Term{canonical}{典范}\Term{equivalence}{等价}于 $\widehat I$。由此得到\Term{topos-morphism}{拓扑斯态射}
\[\operatorname{Top}(m):\operatorname{Top}(X)\longrightarrow\widehat I.\]
设 $F=(F_i)_{i\in I}$（其中 $F_i=F|X_i$）为 $\operatorname{Top}(X)$ 的\Term{object}{对象}（7.4.7）。证明
\[\operatorname{Top}(m)(F)=(i\mapsto\Gamma(X_i,F_i)).\]
由 7.4.12 推出，若 $F$ 为\Term{injective-abelian-sheaf}{内射阿贝尔层}，则每个 $F_i$（$i\in\operatorname{ob}I$）也是\Term{injective-abelian-sheaf}{内射阿贝尔层}。进而推出
\[R^q\operatorname{Top}(m)(F)=(i\mapsto H^q(X_i,F_i)).\]
证明\Term{morphism}{态射} $m$ 的 Cartan--Leray \Term{spectral-sequence}{谱序列}（第 V 讲 5）是
\[H^{p+q}(\operatorname{Top}(X),F)\Longleftarrow\varprojlim_I{}^{(p)}H^q(X_i,F|X_i).\]
\end{SGAresult}
\SGAsection{8}{纤维拓扑斯的投射极限}
\textbf{8.1. 一般性讨论。}
\SourcePage{204}
\begin{SGAresult}{定义}{8.1.1}
设 $\mathcal U$ 为一个\Term{universe}{宇宙}，$p:F\to I$ 为一个 $\mathcal U$-\Term{fibered-topos}{纤维拓扑斯}。一对 $(C,m)$，其中 $C$ 是一个 $\mathcal U$-\Term{topos}{拓扑斯}，$m:C\times I\to F$ 是 $I$ 上 $\mathcal U$-\Term{fibered-topos}{纤维拓扑斯}的一个\Term{cartesian-morphism}{笛卡尔态射}，称为\Term{fibered-topos}{纤维拓扑斯} $F$ 的一个\Term{projective-limit}{投射极限}，如果对每个 $\mathcal U$-\Term{topos}{拓扑斯} $D$，\Term{functor}{函子}
\[\mathscr{H}\!omtop(D,C)\longrightarrow\mathscr{H}\!omtop_{\mathrm{cart}/I}(D\times I,F)\tag{8.1.1.1}\]
是一个\Term{equivalence-of-categories}{范畴等价}；这个\Term{functor}{函子}由\Term{base-change}{基变换}\Term{functor}{函子}
\[\mathscr{H}\!omtop(D,C)\longrightarrow\mathscr{H}\!omtop_I(D\times I,C\times I)\]
与同 $m$ 的复合\Term{functor}{函子}（7.1.8）复合而得。（更“几何”的表述见 8.1.3.3。）
\end{SGAresult}
\textbf{8.1.2.} 采用 8.1.1 的记号。设 $D$ 为一个 $\mathcal U$-\Term{topos}{拓扑斯}，$m_1:D\times I\to F$ 为 $\mathcal U$-\Term{fibered-topos}{纤维拓扑斯}的一个\Term{morphism}{态射}。解释定义 8.1.1，立即得知存在一个\Term{topos-morphism}{拓扑斯态射} $n:D\to C$ 和一个 $I$ 上\Term{fibered-topos}{纤维拓扑斯}\Term{morphism}{态射}的\Term{isomorphism}{同构} $\gamma:m_1\xrightarrow\sim m\circ(n\times\operatorname{id}_I)$，使下图交换：
\[\begin{tikzcd}[column sep=large,row sep=large]
D\times I\arrow[rr,"m_1"]\arrow[d,"m\times\operatorname{id}_I"']&&F\\
C\times I\arrow[urr,"m_1"']&&
\arrow[from=1-1,to=2-3,phantom,"\gamma" description]
\end{tikzcd}\tag{8.1.2.1}\]
此外，若 $(n',\gamma')$ 也满足上述性质，则存在唯一\Term{isomorphism}{同构} $\eta:n\xrightarrow\sim n'$，使
\[\gamma\circ m(n\times\operatorname{id})=\gamma'.\tag{8.1.2.2}\]
因此，若 $(D,m_1)$ 是 $F$ 的一个\Term{projective-limit}{投射极限}，则 $n$ 是 $\mathcal U$-\Term{topos}{拓扑斯}的一个\Term{equivalence-of-categories}{等价}（第 IV 讲 3.4）；再给出使 (8.1.2.1) 交换的\Term{isomorphism}{同构} $\gamma$，便将 $(n,\gamma)$ 确定到\Term{canonical}{典范}\Term{isomorphism}{同构}。当 $I$ \Term{cofiltered-category}{余滤过}时，\Term{projective-limit}{投射极限}的存在性将在 8.2.3 证明。

\textbf{8.1.3.} 将\Term{fibered-topos}{纤维拓扑斯} $F$ 的一个\Term{projective-limit}{投射极限}记为
\[\varprojlim_I\!\operatorname{top}F,\qquad\varprojlim_I\!\operatorname{top}F_i,\qquad\text{或 }\underline F,\]
以 $\mu:(\varprojlim_I\operatorname{top}F)\times I\to F$ 表示\Term{canonical}{典范}\Term{morphism}{态射}，对 $I$ 的每个\Term{object}{对象} $i$，以
\[\mu_i:\underline F\simeq\underline F\times i\longrightarrow F_i\]
表示 $\mu$ 在 $i$ 处\Term{fiber}{纤维}上诱导的\Term{topos-morphism}{拓扑斯态射}。
\SourcePage{205}
选取 $F$ 的一个\Term{bisplitting}{双分裂}（7.1.4），即对每个 $f:i\to j$ 选取\Term{topos-morphism}{拓扑斯态射} $f_\bullet=(f^*,f_*):F_i\to F_j$，使 $f^*$ 是\Term{fibered-structure}{纤维结构}的\Term{inverse-image-functor}{逆像函子}。于是，对 $I$ 中每对可复合\Term{morphism}{态射} $i\xrightarrow f j\xrightarrow g k$，有\Term{topos-morphism}{拓扑斯态射}的\Term{isomorphism}{同构}（7.1.3）
\[c_{f,g}:g_\bullet f_\bullet\simeq(gf)_\bullet.\]
因为 $\mu$ 是\Term{fibered-topos}{纤维拓扑斯}的\Term{cartesian-morphism}{笛卡尔态射}，对 $I$ 的每个 $f:i\to j$，有\Term{topos-morphism}{拓扑斯态射}的\Term{isomorphism}{同构}（7.1.6）
\[b_f:f_\bullet\mu_i\xrightarrow\sim\mu_j,\tag{8.1.3.1}\]
且 $b_f$ 族满足与 (7.1.6.2) 类似的关系。

\textbf{8.1.3.2.} 设 $D$ 为一个 $\mathcal U$-\Term{topos}{拓扑斯}，$m:D\times I\to F$ 为\Term{fibered-topos}{纤维拓扑斯}\Term{morphism}{态射}，$m_i:D=D\times\{i\}\to F_i$ 为 $m$ 在各\Term{fiber}{纤维}上诱导的\Term{morphism}{态射}，且对每个 $f:i\to j$，以 $b'_f:f_\bullet m_i\xrightarrow\sim m_j$ 表示\Term{transition-isomorphism}{转移同构}。由 8.1.2，存在一个\Term{topos-morphism}{拓扑斯态射} $n:D\to\underline F$ 和一个\Term{isomorphism}{同构}族
\[\gamma_i:m_i\xrightarrow\sim\mu_in,\qquad i\in\operatorname{ob}I,\]
使对每个 $f:i\to j$ 有
\[b_f\circ f_\bullet(\gamma_i)=\gamma_j\circ b'_f.\]
而且 $n$ 及 $\gamma_i$（$i\in\operatorname{ob}(I)$）确定到唯一\Term{isomorphism}{同构}（8.1.2）。

反之，若给定一个\Term{topos}{拓扑斯} $\underline F$、\Term{morphism}{态射} $\mu_i:\underline F\to F_i$ 和满足通常相容条件的\Term{isomorphism}{同构} $b_f$（8.1.3.1），且这些数据具有上述\Term{universal-property}{泛性质}，则存在唯一的\Term{fibered-topos}{纤维拓扑斯}\Term{morphism}{态射} $\mu:\underline F\times I\to F$，产生这些 $\mu_i,b_f$；而 $(\underline F,\mu)$ 是 $F$ 的一个\Term{projective-limit}{投射极限}。

\textbf{8.1.3.3.} 设 $\mathcal V$ 为满足 $\mathcal U\in\mathcal V$ 的\Term{universe}{宇宙}。以上讨论允许将\Term{fibered-topos}{纤维拓扑斯}的\Term{projective-limit}{投射极限}解释为如下\Term{pseudofunctor}{伪函子}的 $2$-\Term{category}{范畴}意义的“\Term{projective-limit}{投射极限}”：选取各 $f_\bullet$，从\Term{fibered-structure}{纤维结构}\Term{derivation}{导出}一个取值于属于 $\mathcal V$ 的 $\mathcal U$-\Term{fibered-topos}{纤维拓扑斯}的 $2$-\Term{category}{范畴}的\Term{pseudofunctor}{伪函子}。不过读者须注意，不要将此 $2$-\Term{category}{范畴}意义的“\Term{projective-limit}{投射极限}”（带有交换\Term{isomorphism}{同构}）与严格（或朴素）\Term{projective-limit}{投射极限}（图严格交换）混淆。即使严格概念有意义（\Term{pseudofunctor}{伪函子}为真正\Term{functor}{函子}时），广义“\Term{projective-limit}{投射极限}”与严格\Term{projective-limit}{投射极限}一般也不相同。
\SourcePage{206}
\textbf{8.1.4.} 设 $p:F\to I$、$q:G\to I$ 为两个 $\mathcal U$-\Term{fibered-topos}{纤维拓扑斯}，$m:F\to G$ 为一个\Term{fibered-topos}{纤维拓扑斯}\Term{morphism}{态射}。假设 $F,G$ 都有\Term{projective-limit}{投射极限}（8.1.1）。由定义 8.1.1 立即可知，存在\Term{topos-morphism}{拓扑斯态射}
\[\underline m:\underline F\longrightarrow\underline G,\tag{8.1.4.1}\]
及\Term{fibered-topos}{纤维拓扑斯}\Term{morphism}{态射}的一个\Term{isomorphism}{同构} $\gamma:m\circ\mu\xrightarrow\sim\mu\circ\underline m\times\operatorname{id}_I$，使下图交换：
\[\begin{tikzcd}[column sep=large,row sep=large]
\underline F\times I\arrow[r,"\mu"]\arrow[d,"\underline m\times\operatorname{id}_I"']&F\arrow[d,"m"]\\
\underline G\times I\arrow[r,"\mu"']&G
\arrow[from=1-1,to=2-2,phantom,"\gamma" description]
\end{tikzcd}\tag{8.1.4.2}\]
此外，若 $(\underline m',\gamma')$ 也满足上述性质，则存在唯一的\Term{topos-morphism}{拓扑斯态射}\Term{isomorphism}{同构} $\eta:\underline m\xrightarrow\sim\underline m'$，使
\[\gamma\circ\mu(\eta\times\operatorname{id})=\gamma'.\]
\textbf{8.1.5.} 设 $p:F\to I$ 为属于 $\mathcal V$ 的一个 $\mathcal U$-\Term{fibered-topos}{纤维拓扑斯}，$\phi:I'\to I$ 为属于 $\mathcal V$ 的\Term{functor}{函子}，$p_\phi:F_\phi\to I'$ 为 $(F,p)$ 经\Term{base-change}{基变换} $\phi$ 得到的\Term{fibered-topos}{纤维拓扑斯}（7.1.9）。设 $(\underline F,\mu)$ 为 $F$ 的一个\Term{projective-limit}{投射极限}。经\Term{base-change}{基变换}得到 $I'$ 上的\Term{fibered-topos}{纤维拓扑斯}\Term{morphism}{态射}
\[\mu':\underline F\times I'\longrightarrow F.\]
设 $(\underline F_\phi,\mu_\phi)$ 为 $F_\phi$ 的一个\Term{projective-limit}{投射极限}。由\Term{projective-limit}{投射极限}的\Term{universal-property}{泛性质}（8.1.2），得到\Term{topos-morphism}{拓扑斯态射}
\[m_\phi:\underline F\longrightarrow\underline F_\phi\]
及 $I'$ 上\Term{fibered-topos}{纤维拓扑斯}\Term{morphism}{态射}的一个\Term{isomorphism}{同构}
\[\gamma:\mu'\xrightarrow\sim\mu_\phi\circ(m_\phi\times\operatorname{id}_{I'}).\]
$(m_\phi,\gamma)$ 确定到唯一\Term{isomorphism}{同构}。

\textbf{8.1.6.} 设 $p:F\to I$ 为\Term{small-category}{小范畴} $I$ 上的一个 $\mathcal U$-\Term{fibered-site}{纤维景}，$F^{\sim/I}\to I$ 为相伴 $\mathcal U$-\Term{fibered-topos}{纤维拓扑斯}（7.2.6），$D$ 为一个 $\mathcal U$-\Term{topos}{拓扑斯}。首先有\Term{functor}{函子}
\[\mathscr{H}\!omtop_{\mathrm{cart}/I}(D\times I,F^{\sim/I})\longrightarrow\mathscr{M}orsite_{\mathrm{cart}/I}(D\times I,F),\tag{8.1.6.1}\]
将 $m$ 映到 $m^*$ 与 $\varepsilon_{F/I}$ 的复合\Term{functor}{函子}（7.2.7）。此外，依 $\mathscr{M}orsite_{\mathrm{cart}/I}(D\times I,F)$ 的定义（7.2.2.2），有\Term{functor}{函子}
\[\mathscr{M}orsite_{\mathrm{cart}/I}(D\times I,F)\longrightarrow\mathscr{H}\!om_{\mathrm{cart}/I}(F,D\times I)^\circ.\tag{8.1.6.2}\]
最后，依\Term{inductive-limit}{归纳极限}的定义（6.3），有\Term{functor}{函子}
\SourcePage{207}
\[\mathscr{H}\!om_{\mathrm{cart}/I}(F,D\times I)^\circ\longrightarrow\mathscr{H}\!om(\varinjlim_{I^\circ}F,D)^\circ.\tag{8.1.6.3}\]
以
\[\Theta:\mathscr{H}\!omtop_{\mathrm{cart}/I}(D\times I,F^{\sim/I})\longrightarrow\operatorname{Hom}(\varinjlim_{I^\circ}F,D)^\circ\tag{8.1.6.4}\]
表示这三个\Term{functor}{函子}的复合，以
\[\Pi:F\longrightarrow\varinjlim_{I^\circ}F\tag{8.1.6.5}\]
表示\Term{canonical}{典范}\Term{functor}{函子}。
\begin{SGAresult}{命题}{8.1.7}
\Term{functor}{函子} $\Theta$ \Term{fully-faithful}{完全忠实}。其\Term{essential-image}{本质像}为满足下列条件的全部\Term{functor}{函子} $m:\varinjlim_{I^\circ}F\to D$：\Term{functor}{函子}
\[(m\circ\Pi,p):F\longrightarrow D\times I\]
是从\Term{total-site}{全景} $D\times I$ 到\Term{total-site}{全景} $F$ 的一个\Term{site-morphism}{景态射}。
\end{SGAresult}
\begin{proof}
由 7.2.7，(8.1.6.1) 是\Term{equivalence-of-categories}{范畴等价}；由 7.2.2，(8.1.6.2) \Term{fully-faithful}{完全忠实}；由 6.3，(8.1.6.4) 是\Term{equivalence-of-categories}{范畴等价}。所以 $\Theta$ \Term{fully-faithful}{完全忠实}。$m:\varinjlim_{I^\circ}F\to D$ \Term{isomorphism}{同构}于 $(m\circ\Pi,p)$ 在 $\mathscr{H}\!om_{\mathrm{cart}/I}(F,D\times I)^\circ\to\mathscr{H}\!om(\varinjlim_{I^\circ}F,D)^\circ$ 下的像；而这样的\Term{functor}{函子}属于 (8.1.6.2) 的\Term{essential-image}{本质像}，当且仅当它是从\Term{total-site}{全景} $D\times I$ 到\Term{total-site}{全景} $F$ 的\Term{site-morphism}{景态射}（7.4.13.1）。命题得证。
\end{proof}
\textbf{8.2. 当\Term{index-category}{指标范畴}\Term{cofiltered-category}{余滤过}时\Term{inverse-limit}{投射极限}的构造。}
\begin{SGAresult}{引理}{8.2.1}
设 $I,I'$ 是两个\Term{cofiltered-category}{余滤过范畴}（第 I 讲 8.1），$\phi:I'\to I$ 为满足 $\phi^\circ:I'^\circ\to I^\circ$ \Term{cofinal-functor}{共尾}的\Term{functor}{函子}（第 I 讲 8.1）。又设 $p:F\to I$ 为一个 $\mathcal U$-\Term{fibered-topos}{纤维拓扑斯}，$D$ 为一个 $\mathcal U$-\Term{topos}{拓扑斯}，$m:D\times I\to F$ 为\Term{fibered-topos}{纤维拓扑斯}的一个\Term{cartesian-morphism}{笛卡尔态射}，$m':D\times I'\to F'$ 为 $m$ 经\Term{base-change}{基变换} $\phi$ 得到的\Term{fibered-topos}{纤维拓扑斯}\Term{morphism}{态射}（7.1.9）。$(D,m)$ 是 $F$ 的\Term{projective-limit}{投射极限}，当且仅当 $(D,m')$ 是 $F'$ 的\Term{projective-limit}{投射极限}。
\end{SGAresult}
\begin{proof}
证明只是相当冗长的例行核对，留给有耐心的读者；可采用\Term{pseudofunctor}{伪函子}的\Term{projective-limit}{投射极限}观点（8.1.3.3）。
\end{proof}
\begin{SGAresult}{引理}{8.2.2}
设 $p:F\to I$ 为小\Term{cofiltered-category}{余滤过范畴}上的一个 $\mathcal U$-\Term{fibered-site}{纤维景}，$\underline F$ 为\Term{fibered-category}{纤维范畴} $F$ 的\Term{inductive-limit}{归纳极限}（6.3.9），$\pi:F\to\underline F$ 为\Term{canonical}{典范}\Term{functor}{函子}。在 $\underline F$ 上赋予使 $\pi$ \Term{continuous-functor}{连续}的最粗\Term{topology}{拓扑}（第 III 讲 1）。设 $D$ 为一个 $\mathcal U$-\Term{site}{景}，$n:\underline F\to D$ 为一个\Term{functor}{函子}。下列性质\Term{equivalence}{等价}：

i) $n$ 是\Term{site-morphism}{景态射} $D\to\underline F$。
\SourcePage{208}
ii) $n\circ\pi$ 是\Term{site-morphism}{景态射}（其中 $F$ 配备\Term{total-topology}{全拓扑}）。

iii) $(n\circ\pi,p):F\to D\times I$ 是从\Term{total-site}{全景} $D\times I$ 到\Term{total-site}{全景} $F$ 的\Term{site-morphism}{景态射}。

iv) 对 $I$ 的每个\Term{object}{对象} $i$，复合\Term{functor}{函子} $F_j\xrightarrow{\alpha_{i!}}F\xrightarrow{n\circ\pi}D$ 是\Term{site-morphism}{景态射}。
\end{SGAresult}
\begin{proof}
由 6.6 有 i) $\Leftrightarrow$ ii)，由 7.4.13 有 iii) $\Leftrightarrow$ iv)。证明 iii) $\Leftrightarrow$ ii)。为此只须证明第一\Term{projection}{投影}\Term{functor}{函子} $\operatorname{pr}_1:D\times I\to D$ 是\Term{site-morphism}{景态射}。以 $D^\sim$ 表示 $D$ 上的\Term{sheaf}{层}\Term{topos}{拓扑斯}。由 7.4.7，$(D\times I)^\sim$ \Term{equivalence}{等价}于\Term{topos}{拓扑斯} $\mathscr{H}\!om_{I^\circ}(I^\circ,D^\sim\times I^\circ)=\mathscr{H}\!om(I^\circ,D^\sim)$。此外，立即可知 $\operatorname{pr}_1:D\times I\to D$ 在\Term{sheaf}{层}上的自然延伸（第 III 讲 1.2），是将\Term{functor}{函子} $i\mapsto\sigma(i)\in\operatorname{Hom}(I^\circ,D^\sim)$ 映到\Term{object}{对象} $\varinjlim_{I^\circ}\sigma(i)$ 的\Term{functor}{函子} $\mathscr{H}\!om(I^\circ,D^\sim)\to D^\sim$。由于 $I^\circ$ \Term{filtered-category}{滤过}，沿 $I^\circ$ 的\Term{inductive-limit}{归纳极限}\Term{exact-functor}{正合}（第 II 讲 4），故 $\operatorname{pr}_1$ 是\Term{site-morphism}{景态射}（第 IV 讲 4.9）。

还须证明 ii) $\Leftrightarrow$ iv)。先考虑 $D$ 是一个 $\mathcal U$-\Term{topos}{拓扑斯}，$F$ 是 $I$ 上的一个 $\mathcal U$-\Term{fibered-topos}{纤维拓扑斯}的情形。设 $i$ 为 $I$ 的\Term{object}{对象}。与 7.4.13.6 同理，$\alpha_{i!}:F_i\to F$ \Term{resolution}{分解}为\Term{morphism}{态射} $\overline\alpha_{i!}:F_i\to F_{/\alpha_{i!}(e_i)}$ 和\Term{localization}{局部化}\Term{functor}{函子} $F_{/\alpha_{i!}(e_i)}\to F$，其中 $e_i$ 为 $F_i$ 的一个\Term{final-object}{终对象}。故只须证明复合\Term{functor}{函子} $F_{/\alpha_{i!}(e_i)}\xrightarrow{n\circ\pi}D$ 为\Term{site-morphism}{景态射}。有一个精确到\Term{isomorphism}{同构}交换的图
\[\begin{tikzcd}[column sep=large]
F_{/\alpha_{i!}(e_i)}\arrow[r,"n\circ\pi_{/\alpha_{i!}(e_i)}"]\arrow[d]&D_{/n\circ\pi\alpha_{i!}(e_i)}\arrow[d]\\
F\arrow[r,"n\circ\pi"']&D.
\end{tikzcd}\]
由于是\Term{morphism}{态射}这一性质可以\Term{localization}{局部化}（第 IV 讲 5.10），只须证明 $n\circ\pi\circ\alpha_{i!}(e_i)$ 是 $D$ 的\Term{final-object}{终对象}。对 $I$ 的每个\Term{morphism}{态射} $f:j\to k$，在 $f$ 上方存在唯一\Term{morphism}{态射} $e_f:\alpha_{j!}(e_j)\to\alpha_{k!}(e_k)$，且它是\Term{cartesian-morphism}{笛卡尔态射}。所以 $j\mapsto\alpha_{j!}(e_j)$ 是 $F$ 在 $I$ 上的一个\Term{cartesian-section}{笛卡尔截面}。于是\Term{canonical}{典范}\Term{morphism}{态射}
\[n\circ\pi\circ\alpha_{i!}(e_i)\longrightarrow\varinjlim_I n\circ\pi\circ\alpha_{j!}(e_j)\]
是\Term{isomorphism}{同构}，因为 $\pi$ 将\Term{cartesian-morphism}{笛卡尔态射}映成\Term{isomorphism}{同构}。另一方面，在 $\widehat F$ 中取的\Term{inductive-limit}{归纳极限} $\varinjlim_I\alpha_{j!}(e_j)$ 是 $\widehat F$ 的\Term{final-object}{终对象}。因 $n\circ\pi$ 是\Term{morphism}{态射}，它将 $F^\sim$ 的\Term{final-object}{终对象}映到 $D$ 的\Term{final-object}{终对象}。故 $\varinjlim_I n\circ\pi\circ\alpha_{j!}(e_j)$ 是 $D$ 的\Term{final-object}{终对象}，完成此情形的证明。

一般情形下，\Term{functor}{函子} $(n\pi,p):F\to D\times I$ \Term{cartesian-functor}{笛卡尔}且\Term{continuous-functor}{连续}（7.4.4）。又 $n\circ\pi$ 是复合\Term{functor}{函子}
\[F\xrightarrow{(n\pi,p)}D\times I\xrightarrow{\operatorname{pr}_1}D.\]
\SourcePage{209}
转到相伴 $\mathcal U$-\Term{fibered-topos}{纤维拓扑斯}，得到精确到\Term{isomorphism}{同构}交换的图
\[\begin{tikzcd}[column sep=large]
F\arrow[r,"{(n\pi,p)}"]\arrow[d,"\varepsilon_{F/I}"']&D\times I\arrow[r,"\operatorname{pr}_1"]\arrow[d,"\varepsilon_D\times\operatorname{id}"]&D\arrow[d,"\varepsilon_D"]\\
F^{\sim/I}\arrow[r,"{(n\pi,p)^{\sim/I}}"']&D^\sim\times I\arrow[r,"\operatorname{pr}_1"']&D^\sim.
\end{tikzcd}\]
由于 $(n\pi,p)^{\sim/I}$ \Term{cartesian-functor}{笛卡尔}，它具有形式 $(n'\pi',p')$，其中 $p':F^{\sim/I}\to I$、$\pi':F^{\sim/I}\to\varinjlim_{I^\circ}F_i^\sim$ 是\Term{canonical}{典范}\Term{functor}{函子}。所以 $\operatorname{pr}_1\circ(n\pi,p)^{\sim/I}=n'$。由于 $\varepsilon_{F/I},\varepsilon_D$ 在\Term{sheaf}{层}\Term{topos}{拓扑斯}上诱导\Term{equivalence-of-categories}{等价}（7.4.11），$n,n'$ 在相伴\Term{sheaf}{层}上给出\Term{isomorphism}{同构}的\Term{functor}{函子}（第 III 讲 1.2），故 $n'$ 是\Term{site-morphism}{景态射}。由前述证明，$(n\pi,p)^{\sim/I}$ 在各\Term{fiber}{纤维}上诱导\Term{topos-morphism}{拓扑斯态射}，所以是从\Term{total-site}{全景} $F$ 到\Term{total-site}{全景} $D^\sim\times I$ 的\Term{morphism}{态射}（7.4.13）。由于 $\varepsilon_{F/I},\varepsilon_D\times\operatorname{id}$ 在\Term{total-topos}{全拓扑斯}上诱导\Term{equivalence-of-categories}{等价}（7.4.11），$(n\pi,p)$ 是\Term{fibered-site}{纤维景}\Term{morphism}{态射}（7.4.13），断言得证。
\end{proof}
\begin{SGAresult}{定理}{8.2.3}
设 $p:F\to I$ 为\Term{essentially-small}{本质小}的\Term{cofiltered-category}{余滤过范畴}上的一个 $\mathcal U$-\Term{fibered-site}{纤维景}（第 I 讲 8.1）。

1) 以 $\overrightarrow F$ 表示如下\Term{site}{景}：其底\Term{category}{范畴}为\Term{fibered-category}{纤维范畴} $F\to I$ 的\Term{inductive-limit}{归纳极限}（6.3），其\Term{topology}{拓扑}为使\Term{canonical}{典范}\Term{functor}{函子} $\pi:F\to\overrightarrow F$ \Term{continuous-functor}{连续}的最粗\Term{topology}{拓扑}。则 $\overrightarrow F$ 是一个 $\mathcal U$-\Term{site}{景}，且
\[(\pi,p):F\longrightarrow\overrightarrow F\times I\]
是从\Term{fibered-site}{纤维景} $\overrightarrow F\times I$ 到\Term{fibered-site}{纤维景} $F$ 的一个\Term{morphism}{态射}。

2) 设 $\mu:(\overrightarrow F)^\sim\times I\to F^{\sim/I}$ 是 $(\pi,p)$ 在相伴 $\mathcal U$-\Term{fibered-topos}{纤维拓扑斯}上诱导的\Term{morphism}{态射}。则 $((\overrightarrow F)^\sim,\mu)$ 是 $F^{\sim/I}$ 的一个\Term{projective-limit}{投射极限}。
\end{SGAresult}
\begin{SGAresult}{注记}{8.2.4}
1) 根据 8.1.3 引入的记号，可以置
\[\varprojlim_I\!\operatorname{top}F^\sim=(\overrightarrow F)^\sim.\]
2) 由 7.4.4，$\overrightarrow F$ 上的\Term{topology}{拓扑}也是使所有复合\Term{functor}{函子} $\pi_j:F_j\hookrightarrow F\xrightarrow\pi\overrightarrow F$（$j\in\operatorname{ob}I$）\Term{continuous-functor}{连续}的最粗\Term{topology}{拓扑}。
\end{SGAresult}
\begin{SGAresult}{定义}{8.2.5}
8.2.3 中引入的\Term{site}{景} $\overrightarrow F$ 称为\Term{fibered-site}{纤维景} $F$ 的\Term{projective-limit-site}{投射极限景}。
\end{SGAresult}
作为练习，读者可如同 8.2.1 定义\Term{fibered-site}{纤维景}的\Term{projective-limit}{投射极限}，并验证\Term{site}{景} $\overrightarrow F$ 确实是\Term{site}{景} $F$ 的一个\Term{projective-limit}{投射极限}。
\SourcePage{210}
于是定理的第二个断言可改述为：\Term{projective-limit-site}{投射极限景}上的\Term{sheaf}{层}\Term{topos}{拓扑斯}\Term{equivalence}{等价}于相伴\Term{fibered-topos}{纤维拓扑斯}的\Term{projective-limit}{投射极限}。

\textbf{8.2.6. 定理的证明：归结为小\Term{index-category}{指标范畴}。} 只给出证明的提示。设 $I'$ 为 $I$ 的一个\Term{full-subcategory}{满子范畴}，使 $I'^\circ$ 在 $I^\circ$ 中\Term{cofinal-functor}{共尾}。以 $F'\to I'$ 表示 $F$ 经\Term{base-change}{基变换} $I'\to I$ 得到的\Term{fibered-category}{纤维范畴}，$\pi':F'\to\overrightarrow F'$ 为\Term{canonical}{典范}\Term{functor}{函子}。先可见\Term{category}{范畴} $\overrightarrow F$、$\overrightarrow F'$ \Term{equivalence}{等价}，而且通过此\Term{equivalence-of-categories}{等价}从 $\overrightarrow F$ 的\Term{topology}{拓扑}在 $\overrightarrow F'$ 上诱导的\Term{topology}{拓扑}，就是使 $\pi':F'\to\overrightarrow F'$ \Term{continuous-functor}{连续}的最粗\Term{topology}{拓扑}。所以只要在 $I$ 为小时证明了定理的断言 1)，一般情形即随之得到。一般情形的断言 2) 则由\Term{small-category}{小范畴}情形的断言 2) 和引理 8.2.1 得到。

\textbf{8.2.7. 定理证明的结束。} 第一个断言由引理 8.2.2 取 $D=\overrightarrow F$、$n=\operatorname{id}$ 得到。设 $D$ 为一个 $\mathcal U$-\Term{topos}{拓扑斯}。有精确到\Term{isomorphism}{同构}交换的图
\[\resizebox{\linewidth}{!}{\begin{tikzcd}[ampersand replacement=\&,column sep=small,row sep=large]
\mathscr{H}\!omtop(D,(\overrightarrow F)^\sim)\arrow[r,"(1)"]\arrow[d,"(3)"']\&\mathscr{M}orsite(D,\overrightarrow F)\arrow[r,hook,"(2)"]\arrow[d,"(3)"]\&\mathscr{H}\!om(\overrightarrow F,D)^\circ\arrow[dd,"(5)"]\\
\mathscr{H}\!omtop_{\mathrm{Cart}/I}(D\times I,(\overrightarrow F)^\sim\times I)\arrow[r,"(1)"]\arrow[d,"(4)"']\&\mathscr{M}orsite_{\mathrm{Cart}/I}(D\times I,\overrightarrow F\times I)\arrow[d]\&\\
\mathscr{H}\!omtop_{\mathrm{Cart}/I}(D\times I,F^{\sim/I})\arrow[r,"(1)"']\&\mathscr{M}orsite_{\mathrm{Cart}/I}(D\times I,F)\arrow[r,hook,"(2)"']\&\mathscr{H}\!om_{\mathrm{Cart}/I}(F,D\times I)^\circ,
\end{tikzcd}}\]
其中 (1) 型\Term{functor}{函子}都是\Term{equivalence-of-categories}{等价}（第 IV 讲 4.9.4 及 7.2.7）；(2) 型\Term{functor}{函子}都\Term{fully-faithful}{完全忠实}（第 IV 讲 4.9.1 及 7.2.2）；(3) 型\Term{functor}{函子}是\Term{base-change}{基变换}\Term{functor}{函子}；(4) 型\Term{functor}{函子}分别与 $\mu$、$(\pi,p)$ 复合；(5) 是由\Term{inductive-limit}{归纳极限}的\Term{universal-property}{泛性质}\Term{derivation}{导出}的\Term{equivalence-of-categories}{等价}（6.3）。为证明\Term{functor}{函子}
\[\mathscr{H}\!omtop(D,(\overrightarrow F)^\sim)\xrightarrow{(4)\circ(3)}\mathscr{H}\!omtop_{/I}(D\times I,F^{\sim/I})\]
是\Term{equivalence-of-categories}{等价}，只须证明
\[\mathscr{M}orsite(D,\overrightarrow F)\xrightarrow{(4)\circ(3)}\mathscr{M}orsite_{/I}(D\times I,F)\]
是\Term{equivalence-of-categories}{范畴等价}；亦即只须证明两个\Term{fully-faithful}{完全忠实}\Term{functor}{函子}
\[\mathscr{M}orsite(D,\overrightarrow F)\xrightarrow{(5)\circ(2)}\mathscr{H}\!om_{\mathrm{cart}/I}(F,D\times I)^\circ,\]
\[\mathscr{M}orsite_{/I}(D\times I,F)\xrightarrow{(2)}\mathscr{H}\!om_{\mathrm{cart}/I}(F,D\times I)^\circ\]
\SourcePage{211}
具有相同\Term{essential-image}{本质像}。这立即由引理 8.2.2 得到。

\textbf{8.2.8.} 设 $p:F\to I$ 为\Term{category}{范畴} $I$ 上的一个 $\mathcal U$-\Term{fibered-topos}{纤维拓扑斯}，$(\varprojlim_I\operatorname{top}F_i,\mu)$ 为 $F$ 的一个\Term{projective-limit}{投射极限}。于是有 $I$ 上 $\mathcal U$-\Term{fibered-topos}{纤维拓扑斯}的一个\Term{cartesian-morphism}{笛卡尔态射}
\[\mu:\varprojlim_I\!\operatorname{top}F_i\times I\longrightarrow F.\]
通过\Term{direct-image-functor}{正像函子}转到 $I^\circ$ 上的\Term{fibered-category}{纤维范畴}（7.1.3），得到\Term{cartesian-functor}{笛卡尔函子}
\[\mu_*:\varprojlim_I\!\operatorname{top}F_i\times I^\circ\longrightarrow F',\]
即 $\mu$ 的\Term{direct-image-functor}{正像函子}（7.1.7）。因此由\Term{fibered-category}{纤维范畴}的\Term{projective-limit}{投射极限}定义（6.10），得到\Term{canonical}{典范}\Term{functor}{函子}
\[\Theta:\varprojlim_I\!\operatorname{top}F_i\longrightarrow\varprojlim_{I,f_*}F_i=\mathscr{H}\!om_{\mathrm{cart}/I^\circ}(I^\circ,F').\tag{8.2.8.1}\]
\begin{SGAresult}{定理}{8.2.9}
采用 8.2.8 的记号，假设 $I$ 为\Term{essentially-small}{本质小}的\Term{cofiltered-category}{余滤过范畴}。则\Term{functor}{函子} $\Theta$ 是一个\Term{equivalence-of-categories}{范畴等价}。假设 $I$ 为\Term{small-category}{小范畴}。置 $\varprojlim_I\operatorname{top}F_i=(\overrightarrow F)^\sim$（由 8.2.3、8.2.4，这样做是允许的），并将\Term{category}{范畴} $\mathscr{H}\!om_{I^\circ}(I^\circ,F')$ 解释为\Term{total-site}{全景} $F$ 上的\Term{sheaf}{层}\Term{category}{范畴}（7.4.7）。则复合\Term{functor}{函子}
\[(\overrightarrow F)^\sim=\varprojlim_I\!\operatorname{top}F_i\xrightarrow\Theta\mathscr{H}\!om_{\mathrm{cart}/I^\circ}(I^\circ,F')\hookrightarrow\mathscr{H}\!om_{I^\circ}(I^\circ,F')\simeq F^\sim\]
正是由\Term{functor}{函子} $\pi:F\to\overrightarrow F$ 所定义的\Term{site-morphism}{景态射}的\Term{direct-image-functor}{正像函子}。
\end{SGAresult}
\textbf{8.2.10.} 在证明定理 8.2.9 之前先作几句说明。除了在\Term{index-category}{指标范畴}\Term{cofiltered-category}{余滤过}时解决\Term{projective-limit}{投射极限}的存在性问题（奇怪的是，这在本理论中并非平凡）外，定理 8.2.3、8.2.9 还给出在实际运用\Term{projective-limit}{投射极限}\Term{topos}{拓扑斯}时必不可少的另两项信息。首先，定理 8.2.3 将\Term{projective-limit}{投射极限}\Term{topos}{拓扑斯}解释为某个\Term{site}{景}上的\Term{sheaf}{层}\Term{topos}{拓扑斯}；该\Term{site}{景}的底\Term{category}{范畴}是各 $F_i$ 沿\Term{inverse-image-functor}{逆像函子}系统的\Term{category}{范畴}\Term{inductive-limit}{归纳极限}。在应用中，我们将能具体解释这个\Term{category}{范畴}及其\Term{topology}{拓扑}。另一方面，定理 8.2.9 断言，只要知道\Term{projective-limit}{投射极限}上的一个\Term{sheaf}{层}在各 $F_i$ 中的\Term{direct-image}{正像}系统（它给出\Term{category}{范畴} $(F^{\sim/I})'$ 在 $I^\circ$ 上的一个\Term{cartesian-section}{笛卡尔截面}），便知道这个\Term{sheaf}{层}。此外，反过来，若对每个\Term{object}{对象} $i$ 给定 $F_i$ 上的一个\Term{sheaf}{层} $X_i$，对每个 $f:i\to j$ 给定\Term{isomorphism}{同构} $X_j\simeq f_*X_i$，且这些\Term{isomorphism}{同构}满足相容条件，
\SourcePage{212}
则本质上存在唯一的\Term{projective-limit}{投射极限}上的\Term{sheaf}{层}，其\Term{direct-image}{正像}产生此 $X_i$ 系统。

\textbf{8.2.11.} 8.2.3 的断言 2) 和 8.2.9 的第一个断言可概括为引人注目的公式
\[\varprojlim_{I,f_*}F_i^\sim\simeq\varprojlim_{I,f_\bullet}\!\operatorname{top}F_i^\sim\simeq\left(\varinjlim_{I^\circ,f^*}F_i\right)^\sim.\tag{8.2.11.1}\]
\textbf{8.2.12. 定理 8.2.9 的证明。} 为证明第一个断言，利用 8.2.1 以及关于\Term{fibered-category}{纤维范畴}\Term{projective-limit}{投射极限}的类似引理，可归结为 $I$ 为\Term{small-category}{小范畴}的情形。细节留给读者。以下假设 $I$ 是\Term{small-category}{小范畴}。已知
\[\pi:F\longrightarrow\varinjlim_{I^\circ}F_i\]
是\Term{site-morphism}{景态射}（8.2.2）。它在\Term{topos}{拓扑斯}上定义的\Term{direct-image-functor}{正像函子}
\[\pi_*:\left(\varinjlim_{I^\circ}F_i\right)^\sim\longrightarrow F^\sim\]
\Term{fully-faithful}{完全忠实}，其\Term{essential-image}{本质像}为 $F$ 上将 $F$ 的\Term{cartesian-morphism}{笛卡尔态射}映成\Term{isomorphism}{同构}的全部\Term{sheaf}{层}（6.6）。这些\Term{sheaf}{层}在\Term{equivalence-of-categories}{等价} $F^\sim\simeq\mathscr{H}\!om_{I^\circ}(I^\circ,F')$ 下对应于 $F'$ 在 $I^\circ$ 上的\Term{cartesian-section}{笛卡尔截面}，这是直接的（7.4.7）。此外，对每个 $j\in\operatorname{ob}(I)$，以 $\alpha_{j!}:F_j\to F$ 表示包含\Term{functor}{函子}，它产生三个\Term{functor}{函子} $(\alpha_{j!},\alpha_j^*,\alpha_{j*})$（7.4.3）。依\Term{equivalence-of-categories}{等价} $F^\sim\simeq\mathscr{H}\!om_{I^\circ}(I^\circ,F')$ 的定义（7.4.7），$\varinjlim_{I^\circ}F_i$ 上的\Term{sheaf}{层} $X$ 在\Term{functor}{函子}
\[\left(\varinjlim_{I^\circ}F_i\right)^\sim\xrightarrow{\pi_*}F^\sim\simeq\mathscr{H}\!om_{I^\circ}(I^\circ,F')\]
下对应的\Term{cartesian-section}{笛卡尔截面}是 $j\mapsto\alpha_j^*\pi_*(X)$。而由 8.2.3 中对\Term{functor}{函子} $\mu:(\varinjlim_{I^\circ}F_i)^\sim\times I\to F$ 的描述可知，对于每个\Term{object}{对象} $j$，\Term{functor}{函子} $\mu_{j*}:(\varinjlim_{I^\circ}F_i)^\sim\to F_j$，即由 $\mu$ 转到 $j$ 处\Term{fiber}{纤维}所得的\Term{morphism}{态射} $\mu_j:(\varinjlim_{I^\circ}F_i)^\sim\to F_j$ 的\Term{direct-image-functor}{正像函子}，正是 $\alpha_j^*\pi_*$。所以\Term{functor}{函子}
\[\left(\varinjlim_{I^\circ}F_i\right)^\sim\xrightarrow{\pi_*}F^\sim\simeq\mathscr{H}\!om_{I^\circ}(I^\circ,F')\]
正是\Term{functor}{函子}
\[\left(\varinjlim_{I^\circ}F_i\right)^\sim\xrightarrow\Theta\mathscr{H}\!om_{\mathrm{cart}/I^\circ}(I^\circ,F')\hookrightarrow\mathscr{H}\!om_{I^\circ}(I^\circ,F'),\]
定理得证。
\textbf{8.3. \Term{projective-limit-site}{投射极限景}的\Term{topology}{拓扑}：\Term{coherent-topos}{凝聚拓扑斯}的情形。}

\textbf{8.3.1.} 本节中，$p:F\to I$ 是\Term{essentially-small}{本质小}的\Term{cofiltered-category}{余滤过范畴}上的一个 $\mathcal U$-\Term{fibered-site}{纤维景}，其\Term{inverse-image-functor}{逆像函子} $f^*:F_i\to F_j$ 与\Term{fibre-product}{纤维积}交换。此外，对每个\Term{fiber}{纤维}\Term{category}{范畴} $F_i$ 的每个\Term{object}{对象} $X$，给定一个由 $F_i$ 的\Term{topology}{拓扑}的\Term{covering-family}{覆盖族}组成的\Term{set}{集合} $\operatorname{Cov}(X)$；这些族具有第 II 讲 1.3 的性质 (PT0)、(PT1)，且生成 $F_i$ 的\Term{topology}{拓扑}。最后，假设对每个 $f:i\to j$、每个 $Y\in\operatorname{ob}F_j$ 和每个族 $(Y_\alpha\to Y)_{\alpha\in A}\in\operatorname{Cov}(Y)$，族 $(f^*(Y_*)\to f^*(Y))_{\alpha\in A}$ 属于 $\operatorname{Cov}(f^*(Y))$。例如，取 $\operatorname{Cov}(X)=\{\text{$F_i$ 中 $X$ 的所有覆盖族}\}$ 时，这些条件成立。
\SourcePage{213}
本节拟在适当的\Term{finite}{有限}性假设下（8.3.13）研究 $F$ 的\Term{projective-limit-site}{投射极限景}。

\textbf{8.3.2.} 设 $\pi:F\to\overrightarrow F$ 为\Term{canonical}{典范}\Term{functor}{函子}，$Y$ 为 $\overrightarrow F$ 的一个\Term{object}{对象}。以 $\operatorname{Cov}(Y)$ 表示如下类型的族 $(Y_\alpha\xrightarrow{n_\alpha}Y)_{\alpha\in A}$ 的\Term{set}{集合}：存在 $i\in I$、$F_i$ 的\Term{object}{对象} $X$、族 $(X_\alpha\xrightarrow{n'_\alpha}X)_{\alpha\in A}\in\operatorname{Cov}(X)$、\Term{isomorphism}{同构} $\phi:\pi(X)\simeq Y$ 和\Term{isomorphism}{同构}族 $\phi_\alpha:\pi(X_\alpha)\simeq Y_\alpha$（$\alpha\in A$），使每个 $\alpha$ 都有\Term{commutative-diagram}{交换图}
\[\begin{tikzcd}
\pi(X_\alpha)\arrow[r,"\phi_\alpha","\sim"']\arrow[d,"\pi(n'_\alpha)"']&Y_\alpha\arrow[d,"n_\alpha"]\\
\pi(X)\arrow[r,"\phi","\sim"']&Y.
\end{tikzcd}\]
\begin{SGAresult}{命题}{8.3.3}
1) 属于各 $\operatorname{Cov}(Y)$（$Y\in\operatorname{ob}\overrightarrow F$）的族的全体生成\Term{projective-limit-site}{投射极限景}的\Term{topology}{拓扑}（8.2.5）。

2) $Y\mapsto\operatorname{Cov}(Y)$（$Y\in\operatorname{ob}\overrightarrow F$）具有第 II 讲 1.3 的性质 (PT0)、(PT1)。
\end{SGAresult}
下面需要一个引理。
\begin{SGAresult}{引理}{8.3.4}
设 $i$ 为 $I$ 的\Term{object}{对象}，$n:X'\to X$ 为 $F_i$ 的一个\Term{quarrable-morphism}{可作基变换的态射}（morphisme quarrable），且每个 $f:j\to i$ 的 $f^*(n)$ 仍\Term{quarrable-morphism}{可作基变换}。则 $\pi(n)$ \Term{quarrable-morphism}{可作基变换}。\Term{functor}{函子}
\[\pi_j:F_j\hookrightarrow F\xrightarrow\pi\overrightarrow F,\qquad j\in\operatorname{ob}I,\]
与\Term{fibre-product}{纤维积}交换。
\end{SGAresult}
\begin{proof}
回顾 $\overrightarrow F$ 的每个\Term{object}{对象} $Y$ 都等于某个\Term{object}{对象} $\pi(Z)$，$Z\in\operatorname{ob}F$，而每个\Term{morphism}{态射} $m:Y\to\pi(X)$ 具有形式
\[\pi(Z)\xrightarrow{\pi(s)^{-1}}\pi(Z')\xrightarrow{\pi(m')}\pi(X),\]
其中 $s$ 为 $F$ 的\Term{cartesian-morphism}{笛卡尔态射}。所以，为证明 $\pi(n)$ \Term{quarrable-morphism}{可作基变换}，只须证明每个图
\[\begin{tikzcd}
&\pi(X')\arrow[d,"\pi(u)"]\\
\pi(Z')\arrow[r,"\pi(m')"']&\pi(X)
\end{tikzcd}\]
的\Term{fibre-product}{纤维积}\Term{representable}{可表}。
\SourcePage{214}
但 $m':Z'\to X$ \Term{decomposable}{可分解}为 $Z'\xrightarrow{m''}X''\xrightarrow{s'}X$，其中 $p(m'')$ 为 $p(Z')=j$ 的恒等，而 $s'$ 是位于 $p(s')=f$ 上方的\Term{cartesian-morphism}{笛卡尔态射}。置 $X'''=f^*(X')$、$n'=f^*(n)$，有\Term{commutative-diagram}{交换图}
\[\begin{tikzcd}
X'''\arrow[r,"s''"]\arrow[d,"n'"']&X'\arrow[d,"n"]\\
X''\arrow[r,"s'"']&X,
\end{tikzcd}\]
从而有 $\overrightarrow F$ 中的\Term{commutative-diagram}{交换图}
\[\begin{tikzcd}
&\pi(X''')\arrow[r,"\sim"]\arrow[d,"\pi(u')"']&\pi(X')\arrow[d,"\pi(n)"]\\
\pi(Z')\arrow[r,"\pi(m')"']&\pi(X'')\arrow[r,"\sim"']&\pi(X),
\end{tikzcd}\]
其中 $n'$ 是 $F_j$ 的一个\Term{quarrable-morphism}{可作基变换的态射}。为证明 $\pi(n)$ \Term{quarrable-morphism}{可作基变换}，只须证明 $\pi(Z')\times_{\pi(X'')}\pi(X''')$ \Term{representable}{可表}，因而只须证明 $\pi_j:F_j\hookrightarrow F\to\overrightarrow F$ 与\Term{fibre-product}{纤维积}交换。

设
\[\begin{tikzcd}
Y'\arrow[r]\arrow[d]&X'\arrow[d]\\
Y\arrow[r]&X
\end{tikzcd}\]
为 $F_j$ 的一个\Term{cartesian-diagram}{笛卡尔图}，$W$ 为 $F$ 中位于 $k\in\operatorname{ob}I$ 上方的\Term{object}{对象}。由 6.5 有
\[\operatorname{Hom}(\pi(W),\pi(Y'))\simeq\varinjlim_{(\ell,\,f:\ell\to j,\,g:\ell\to k)}\operatorname{Hom}_{F_\ell}(g^*(W),f^*(Y')).\]
由于 $f^*:F_j\to F_\ell$ 与\Term{fibre-product}{纤维积}交换，有
\[\begin{aligned}
\operatorname{Hom}(g^*(W),f^*(Y'))\simeq{}&\operatorname{Hom}(g^*(W),f^*(Y))\\
&\times_{\operatorname{Hom}(g^*(W),f^*(X))}\operatorname{Hom}(g^*(W),f^*(X')).
\end{aligned}\]
利用\Term{filtered-limit}{滤过极限}与\Term{fibre-product}{纤维积}的交换性（第 I 讲 2）及公式 6.5，得到\Term{isomorphism}{同构}
\[\begin{aligned}
\operatorname{Hom}(\pi(W),\pi(Y'))\simeq{}&\operatorname{Hom}(\pi(W),\pi(Y))\\
&\times_{\operatorname{Hom}(\pi(W),\pi(X))}\operatorname{Hom}(\pi(W),\pi(X')).
\end{aligned}\]
\SourcePage{215}
这表明图
\[\begin{tikzcd}
\pi(Y')\arrow[r]\arrow[d]&\pi(X')\arrow[d]\\
\pi(Y)\arrow[r]&\pi(X)
\end{tikzcd}\]
在 $\overrightarrow F$ 中为\Term{cartesian-diagram}{笛卡尔图}。
\end{proof}
\textbf{8.3.5. 8.3.3 的证明。} 先证明第二个断言。由 8.3.4，各 $\operatorname{Cov}(Y)$（$Y\in\overrightarrow F$）中的族由\Term{quarrable-morphism}{可作基变换的态射}组成，故有性质 (PT0)。为证明这些族有性质 (PT1)（\Term{base-change}{基变换}下的\Term{stable}{稳定}性），作如同 8.3.3 的证明中的一系列化简，可归结为证明下列断言：对每个 $i\in\operatorname{ob}I$、每个 $Y\in\operatorname{ob}F_i$、每个 $(Y_\alpha\to Y)_{\alpha\in A}\in\operatorname{Cov}(Y)$ 和每个 $F_i$ 中的\Term{morphism}{态射} $u:Z\to Y$，族
\[(\pi(Z)\times_{\pi(Y)}\pi(Y)\longrightarrow\pi(Z))_{\alpha\in A}\]
属于 $\operatorname{Cov}((Z))$。

此断言来自 $\pi_i:F_i\hookrightarrow F\xrightarrow\pi\overrightarrow F$ 与\Term{fibre-product}{纤维积}交换（8.3.3），以及 $F_i$ 中的 $X\mapsto\operatorname{Cov}(X)$ 在\Term{base-change}{基变换}下\Term{stable}{稳定}。证明第一个断言。显然，属于各 $\operatorname{Cov}(Y)$（$Y\in\overrightarrow F$）的族，对于使 $\pi$ \Term{continuous-functor}{连续}的最粗\Term{topology}{拓扑} $T$ 是\Term{covering-family}{覆盖族}（第 III 讲 1）。故由这些族生成的\Term{topology}{拓扑} $T'$ 不细于 $T$。为了证明它细于 $T$（因而等于 $T$），只须证明每个 $T'$-\Term{sheaf}{层} $M$ 都使 $M\circ\pi$ 为 $F$ 上的\Term{sheaf}{层}，亦即（8.2.4）每个 $M\circ\pi_i$ 都是 $F_i$ 上的\Term{sheaf}{层}。由于各 $\pi_i:F_i\to\overrightarrow F$ 与\Term{fibre-product}{纤维积}交换（8.3.4），$M$ 是 $T'$ 的\Term{sheaf}{层}，当且仅当（第 II 讲 2.3 和第 I 讲 2.12）对每个 $i\in\operatorname{ob}I$、每个 $Y\in\operatorname{ob}F_i$、每个 $(Y_\alpha\to Y)_{\alpha\in A}\in\operatorname{Cov}(Y)$，\Term{set}{集合}序列
\[M(\pi_i(Y))\longrightarrow\prod_{\alpha\in A}M(\pi_i(Y_\alpha))\rightrightarrows\prod_{(\alpha,\beta)\in A\times A}M(\pi_i(Y_\alpha\times_Y Y_\beta))\]
\Term{exact-sequence}{正合}，即（同上）当且仅当每个 $M\circ\pi_i$ 都是 $F_i$ 上的\Term{sheaf}{层}。

\begin{SGAresult}{命题}{8.3.6}
采用 8.3.1、8.3.2 的记号和假设。若对每个 $i\in\operatorname{ob}I$，$X\mapsto\operatorname{Cov}(X)$ 是 $F_i$ 上的一个\Term{pretopology}{预拓扑}（第 II 讲 1.3），且对每个 $X\in\operatorname{ob}F_i$，$\operatorname{Cov}(X)$ 的\Term{covering-family}{覆盖族}都是\Term{finite}{有限}族，则对每个 $Y\in\overrightarrow F$，$\operatorname{Cov}(Y)$ 的\Term{covering-family}{覆盖族}都是\Term{finite}{有限}族，而且 $Y\mapsto\operatorname{Cov}(Y)$ 是 $\overrightarrow F$ 上的一个\Term{pretopology}{预拓扑}。
\end{SGAresult}
\textbf{8.3.7.} 唯一须证明的是 $Y\mapsto\operatorname{Cov}(Y)$ 的族有第 II 讲 1.3 的性质 (PT2)（复合下的\Term{stable}{稳定}性）。设 $i\in\operatorname{ob}I$、$Y\in\operatorname{ob}F_i$，且 $(Y_\alpha\to Y)_{\alpha\in A}\in\operatorname{Cov}(Y)$。再设 $i_\alpha$（$\alpha\in A$）为 $I$ 的\Term{object}{对象}族，对于每个 $\alpha$，$Z_\alpha$ 为 $F_i$ 的\Term{object}{对象}，而 $(Z_{\alpha\beta}\xrightarrow{n_{\alpha\beta}}Z_\alpha)_{\alpha\beta\in B_\alpha}$ 为 $\operatorname{Cov}(Z_\alpha)$ 的一个族。
\SourcePage{216}
最后，对每个 $\alpha\in A$ 给定一个\Term{isomorphism}{同构} $\phi_\alpha:\pi(Z_\alpha)\xrightarrow\sim\pi(Y_\alpha)$。回顾各 $\operatorname{Cov}(X)$（$X\in\operatorname{ob}\overrightarrow F$）中族的定义（8.3.2），可知须证明族
\[\left(\pi(Z_{\alpha\beta})\xrightarrow{\pi(n_\alpha)\circ\phi_\alpha\circ\pi(n_{\alpha\beta})}\pi(Y)\right)_{\alpha\beta\in\coprod_A B_\alpha}\]
属于 $\operatorname{Cov}(\pi(Y))$。下面经五次化简证明。

\textbf{8.3.8. 归结为对每个 $\alpha\in A$，$\phi_\alpha=\pi(m_\alpha)$，其中 $m_\alpha:Z_\alpha\to Y_\alpha$。} 对每个 $\alpha$ 有 $\phi_\alpha=\pi(m_\alpha)\pi(s_\alpha)^{-1}$（6.5），其中 $s_\alpha$ 为位于 $f_\alpha:i'_\alpha\to i_\alpha$ 上方的\Term{cartesian-morphism}{笛卡尔态射}。故有\Term{commutative-diagram}{交换图}
\[\begin{tikzcd}
Z_{\alpha\beta}\arrow[d,"n_{\alpha\beta}"']&f_\alpha^*(Z_{\alpha\beta})\arrow[l,"s_{\alpha\beta}"']\arrow[d,"f_\alpha^*(n_{\alpha\beta})"]\\
Z_\alpha&f_\alpha^*(Z_\alpha)\arrow[l,"s_\alpha"]
\end{tikzcd}\qquad\alpha\beta\in\coprod_A B_\alpha,\]
从而有\Term{commutative-diagram}{交换图}
\[\begin{tikzcd}[column sep=large]
\pi(Z_{\alpha\beta})\arrow[r,"\sim"]\arrow[d,"\pi(n_{\alpha\beta})"']&\pi(f_\alpha^*(Z_{\alpha\beta}))\arrow[d,"\pi(f_\alpha^*(m_{\alpha\beta}))"]&\\
\pi(Z_\alpha)\arrow[r,"\pi(s_\alpha)^{-1}","\sim"']&\pi(f_\alpha^*(Z_\alpha))\arrow[r,"\pi(m_\alpha)"']&\pi(Y_\alpha),
\end{tikzcd}\]
其中 $\alpha\beta\in\coprod_A B_\alpha$。由于族 $(f_\alpha^*(Z_{\alpha\beta})\xrightarrow{f_\alpha^*(n_\alpha)}f^*(Z_\alpha))$ 属于 $\operatorname{Cov}(f_\alpha^*(Z_\alpha))$，可归结为对每个 $\alpha$ 有 $\phi_\alpha=\pi(m_\alpha)$ 的情形。

\textbf{8.3.9. 归结为对每个 $\alpha\in A$，$i_\alpha=j$ 且 $p(m_\alpha)=k:j\to i$。} 置 $g_\alpha=p(m_\alpha)$。由于 $I$ \Term{cofiltered-category}{余滤过}而 $A$ \Term{finite}{有限}，存在 $I$ 的\Term{object}{对象} $j$ 和\Term{morphism}{态射} $h_\alpha:j\to i_\alpha$，使每对 $(\alpha,\alpha')$ 都有 $g_\alpha h_\alpha=g_{\alpha'}h_{\alpha'}=k$。如同 8.3.8，利用\Term{inverse-image-functor}{逆像函子} $h_\alpha^*$ 即可归结为所述情形。

\textbf{8.3.10. 归结为 $i=j$ 且 $k:j\to i$ 为恒等。} 对每个 $\alpha$，有\Term{commutative-diagram}{交换图}
\[\begin{tikzcd}
Z_\alpha\arrow[r,"m'_\alpha"]&k^*(Y_\alpha)\arrow[r,"t_\alpha"]\arrow[d,"k^*(n_\alpha)"']&Y_\alpha\arrow[d,"n_\alpha"]\\
&k^*(Y)\arrow[r,"t"']&Y,
\end{tikzcd}\]
\SourcePage{217}
其中 $t,t_\alpha$（$\alpha\in A$）是\Term{cartesian-morphism}{笛卡尔态射}，$m'_\alpha$ 位于 $j$ 的恒等上方，且 $t_\alpha\circ m'_\alpha=m_\alpha$。用族 $(k^*(Y_\alpha)\to k^*(Y))_{\alpha\in A}\in\operatorname{Cov}(k^*(Y_\alpha))$ 代替 $(Y_\alpha\to Y)_{\alpha\in A}$，用 $m'_\alpha$ 代替 $m_\alpha$，便归结为所述情形。

\begin{SGAresult}{引理}{8.3.11}
设 $j$ 为 $I$ 的\Term{object}{对象}，$m:X'\to X$ 为 $F_j$ 的\Term{morphism}{态射}，且 $\pi(m)$ 是\Term{isomorphism}{同构}。则存在\Term{morphism}{态射} $\ell:j'\to j$，使 $\ell^*(m)$ 为\Term{isomorphism}{同构}。
\end{SGAresult}
\begin{proof}
先证明存在\Term{morphism}{态射} $f:i\to j$ 和\Term{morphism}{态射} $q:f^*(X)\to f^*(X')$，使 $f^*(m)q=\operatorname{id}_{f^*(X)}$ 且 $\pi(q)$ 是\Term{isomorphism}{同构}。事实上，存在\Term{isomorphism}{同构} $n:\pi(X)\to\pi(X')$ 满足 $\pi(m)\circ n=\operatorname{id}_{\pi(X)}$。所以（6.5）存在两个\Term{morphism}{态射} $X\xleftarrow{s_1}Y_1\xrightarrow{q_1}X'$，其中 $s_1$ \Term{cartesian-morphism}{笛卡尔}，使 $n=\pi(q_1)\pi(s_1)^{-1}$。于是 $\pi(m)\circ\pi(q_1)=\pi(s_1)$。从 $Y_1$ 到 $X$ 的\Term{morphism}{态射} $mq_1,s_1$ 在 $\overrightarrow F$ 中有相同的像。因此（参见 6.5 中\Term{inductive-limit}{归纳极限}\Term{morphism}{态射}的描述），存在\Term{cartesian-morphism}{笛卡尔态射} $s_2:Y_2\to Y_1$，使 $mq_1s_2=s_1s_2$。置 $q_1s_2=q_2$，以 $s_3$ 表示\Term{cartesian-morphism}{笛卡尔态射} $s_1s_2$，则 $mq_2=s_3$。$q_2:Y_2\to X'$ 本质唯一地\Term{resolution}{分解}为 $Y_2\xrightarrow qY_3\xrightarrow{s_4}X'$，其中 $s_4$ \Term{cartesian-morphism}{笛卡尔}，$q$ 位于 $p(Y_2)$ 的恒等上方。故有\Term{commutative-diagram}{交换图}
\[\begin{tikzcd}
Y_3\arrow[d,"s_4"']&Y_2\arrow[l,"q"']\arrow[d,"s_3"]\\
X'\arrow[r,"m"']&X.
\end{tikzcd}\tag{8.3.11.1}\]
置 $p(s_4)=p(s_3)=f:i\to j$。有 $f^*(X')\simeq Y_3$、$f^*(X')\simeq Y_2$，而 $f^*(m):Y_2\to Y_3$ 使下图交换：
\[\begin{tikzcd}
Y_3\arrow[r,"f^*(m)"]\arrow[d,"s_4"']&Y_2\arrow[d,"s_3"]\\
X'\arrow[r,"m"']&X.
\end{tikzcd}\]
所以 $s_3=s_3(f^*(m)\circ q)$；因 $s_3$ \Term{cartesian-morphism}{笛卡尔}且 $f^*(m)\circ q$ 位于 $p(Y_2)$ 的恒等上方，有 $f^*(m)\circ q=\operatorname{id}_{Y_2}$。此外，由 (8.3.11.1) 的交换性，$\pi(q)$ 是\Term{isomorphism}{同构}。

将刚才的结果应用于 $q:f^*(X)\to f^*(X')$，得到\Term{morphism}{态射} $f':j'\to i$ 和\Term{morphism}{态射} $q':f'^*(X')\to f'^*(X)$，使 $f'^*(q)q'=\operatorname{id}_{f'^*(f'^*(X'))}$。置 $\ell^*=ff'$，有\Term{isomorphism}{同构} $f'^*f^*\simeq\ell^*$。所以有\Term{morphism}{态射} $f'^*(q):\ell^*(X)\to\ell^*(X')$ 和\Term{isomorphism}{同构} $f'^*f^*\simeq\ell^*$。于是
\SourcePage{218}
有\Term{morphism}{态射} $f'^*(q):\ell^*(X)\to\ell^*(X')$ 和 $q':\ell^*(X')\to\ell^*(X)$，使
\[\ell^*(m)f'^*(q)=\operatorname{id}_{\ell^*(X)},\qquad f'^*(q)q'=\operatorname{id}_{\ell^*(X')}.\]
所以 $\ell^*(m)$ 是\Term{isomorphism}{同构}。
\end{proof}
\textbf{8.3.12. 证明的结束。} 由于 $I$ \Term{cofiltered-category}{余滤过}而 $A$ \Term{finite}{有限}，由 8.3.11 存在一个\Term{inverse-image-functor}{逆像函子} $\ell^*$ 将所有 $m_\alpha$（$\alpha\in A$）都映成\Term{isomorphism}{同构}。故可归结为各 $m_\alpha$ 为\Term{isomorphism}{同构}的情形。此时\Term{covering-family}{覆盖族} $(Z_{\alpha\beta}\to Z)_{\alpha\beta\in B_\alpha}$ 是由\Term{covering-family}{覆盖族} $(Y_{\alpha\beta}\to Y_\alpha)_{\alpha\beta\in B_\alpha}\in\operatorname{Cov}(Y_\alpha)$ 经\Term{base-change}{基变换} $m_\alpha:Z_\alpha\to Y_\alpha$ 得到的。因此可归结为 $Z_\alpha=Y_\alpha$ 且 $\phi_\alpha=\operatorname{id}_{\pi(Y_\alpha)}$ 的情形。但此时族
\[\left(Y_{\alpha\beta}\xrightarrow{n_\alpha\circ n_{\alpha\beta}}Y\right)_{\alpha\beta\in\coprod_A B_\alpha}\]
属于 $\operatorname{Cov}(Y)$（性质 (PT2)），所以它在 $\pi$ 下的像属于 $\operatorname{Cov}(\pi(Y))$（8.3.2）。

\begin{SGAresult}{定理}{8.3.13}
设 $p:F\to I$ 为\Term{essentially-small}{本质小}的\Term{cofiltered-category}{余滤过范畴}上的一个 $\mathcal U$-\Term{fibered-topos}{纤维拓扑斯}。若对每个 $i\in\operatorname{ob}I$，\Term{fiber}{纤维}\Term{topos}{拓扑斯} $F_i$ \Term{coherent-topos}{凝聚}（2.3），且对每个 $f:i\to j$，\Term{topos-morphism}{拓扑斯态射} $f_\bullet=(f^*,f_*):F_i\to F_j$ \Term{coherent-morphism}{凝聚}（3.1），则\Term{topos}{拓扑斯} $\varprojlim_I\operatorname{top}F_i$ \Term{coherent-topos}{凝聚}，且对每个 $i\in\operatorname{ob}I$，\Term{canonical}{典范}\Term{topos-morphism}{拓扑斯态射}
\[\mu_i:\varprojlim_I\!\operatorname{top}F_i\longrightarrow F_i\]
\Term{coherent-morphism}{凝聚}。

此外，以 $F_{\mathrm{coh}}$ 表示由 $F$ 中在其\Term{fiber}{纤维}内\Term{coherent-object}{凝聚}的\Term{object}{对象}（1.13）定义的\Term{full-subcategory}{满子范畴}，则 $F_{\mathrm{coh}}$ 在 $I$ 上\Term{fibered-category}{纤维化}，且 $\overrightarrow F_{\mathrm{coh}}$ \Term{canonical}{典范}\Term{equivalence}{等价}于 $\varprojlim_I\operatorname{top}F_i$ 的\Term{coherent-object}{凝聚对象}\Term{category}{范畴}。
\end{SGAresult}
\begin{proof}
对每个 $i\in\operatorname{ob}I$，以 $(F_i)_{\mathrm{coh}}$ 表示 $F_i$ 的\Term{coherent-object}{凝聚对象}组成的 $\mathcal U$-\Term{site}{景}（赋予\Term{induced-topology}{诱导拓扑}）（1.13）。因 $F_i$ \Term{coherent-topos}{凝聚}，它是 $(F_i)_{\mathrm{coh}}$ 上的\Term{sheaf}{层}\Term{topos}{拓扑斯}。对每个 $f:i\to j$，$f_\bullet=(f^*,f_*)$ \Term{coherent-morphism}{凝聚}，故 $f^*(F_j)_{\mathrm{coh}}\subset(F_i)_{\mathrm{coh}}$。于是有 $\mathcal U$-\Term{fibered-site}{纤维景} $p_{\mathrm{coh}}:F_{\mathrm{coh}}\to I$，其\Term{fiber}{纤维}\Term{site}{景}是 $(F_i)_{\mathrm{coh}}$，相伴\Term{fibered-topos}{纤维拓扑斯}\Term{equivalence}{等价}于 $p:F\to I$。所以（8.2.3）
\[\varprojlim_{I,f_\bullet}\!\operatorname{top}F_i\simeq(\overrightarrow F_{\mathrm{coh}})^\sim.\tag{8.3.13.1}\]
由于各 $F_i$ \Term{coherent-topos}{凝聚}，$(F_i)_{\mathrm{coh}}$ 中\Term{finite-product}{有限积}与\Term{fibre-product}{纤维积}\Term{representable}{可表}（2.2）。又 $f^*:(F_i)_{\mathrm{coh}}\to(F_j)_{\mathrm{coh}}$ \Term{left-exact}{左正合}。最后，由于 $(F_i)_{\mathrm{coh}}$ 的\Term{object}{对象}\Term{quasi-compact}{拟紧}，\Term{finite}{有限}\Term{covering-family}{覆盖族}构成一个\Term{pretopology}{预拓扑}。由 8.3.4，$\varprojlim_{I^\circ,F}(F_i)_{\mathrm{coh}}$ 中\Term{finite-product}{有限积}与\Term{fibre-product}{纤维积}\Term{representable}{可表}（关于引理 8.3.4 的一点补憾：证明\Term{canonical}{典范}\Term{functor}{函子} $(F_i)_{\mathrm{coh}}\to\overrightarrow F_{\mathrm{coh}}$ 将\Term{final-object}{终对象}映到\Term{final-object}{终对象}）；由 8.3.6，$\overrightarrow F_{\mathrm{coh}}$ 的\Term{object}{对象}\Term{quasi-compact}{拟紧}。所以（2.4.5）$\varprojlim_I\operatorname{top}(F_i)$ \Term{coherent-topos}{凝聚}。最后，\Term{canonical}{典范}\Term{morphism}{态射} $\mu_i:\varprojlim_I\operatorname{top}(F_i)\to F_i$
\SourcePage{219}
是由\Term{site-morphism}{景态射}（8.2.3）
\[(F_i)_{\mathrm{coh}}\longrightarrow\overrightarrow F_{\mathrm{coh}}\]
转到\Term{sheaf}{层}\Term{topos}{拓扑斯}得到的，故 $\mu_i$ \Term{coherent-morphism}{凝聚}（3.3）。

证明最后的断言。由 (8.3.13.1)，有
\[\varprojlim_I\!\operatorname{top}(F_i)\simeq(\overrightarrow F_{\mathrm{coh}})^\sim.\]
设 $X$ 是 $(F_{\mathrm{coh}})^\sim$ 的一个\Term{coherent-object}{凝聚对象}。因 $X$ \Term{quasi-compact}{拟紧}，存在\Term{finite}{有限}\Term{covering-family}{覆盖族} $(Y_\alpha\to X)$，其中 $Y_\alpha\in\operatorname{ob}\overrightarrow F_{\mathrm{coh}}$（1.1）。取足够小的指标 $i\in\operatorname{ob}I$，可假设 $Y_\alpha$ 为 $F_i^{\mathrm{coh}}$ 中某个族 $Y'_\alpha$ 在 $\mu_i:F_i^{\mathrm{coh}}\to\overrightarrow F_{\mathrm{coh}}$ 下的像。取各 $Y'_\alpha$ 的\Term{direct-sum}{直和}（1.15），可知存在指标 $i\in\operatorname{ob}I$、$Y'\in F_i^{\mathrm{coh}}$ 和一个满的\Term{morphism}{态射} $\mu_i(Y')\to X$。\Term{equivalence-relation}{等价关系} $R=\mu_i(Y')\times_X\mu_i(Y')$ 为\Term{coherent-object}{凝聚对象}，因为 $(\overrightarrow F_{\mathrm{coh}})^\sim$ 是\Term{coherent-topos}{凝聚拓扑斯}（2.2）。先证明 $R$ 是 $\overrightarrow F_{\mathrm{coh}}$ 的\Term{object}{对象}。依定义，$R$ 是 $\mu_i(Y'\times Y')$ 的\Term{subobject}{子对象}；由于 $R$ \Term{coherent-object}{凝聚}，可以改变指标 $i$，找到 $F_i^{\mathrm{coh}}$ 中的一支\Term{arrow}{箭头} $m:Z\to(Y'\times Y')$，使 $\operatorname{Im}(\mu_i(m))=R$。由于 $F_i$ \Term{coherent-topos}{凝聚}，$\operatorname{Im}(m)$ 在 $F_i$ 中\Term{coherent-object}{凝聚}（1.17.1），因而属于 $F_i^{\mathrm{coh}}$。因此存在指标 $i$ 及 $F_i^{\mathrm{coh}}$ 的一个图 $R\mathrel{\substack{\xrightarrow{p_1}\\[-.7ex]\xrightarrow{p_2}}}Y'$，使
\[\mu_i(R)\mathrel{\substack{\xrightarrow{\mu_i(p_1)}\\[-.7ex]\xrightarrow{\mu_i(p_2)}}}\mu_i(Y')\]
为\Term{equivalence-relation}{等价关系}，且 $X=\mu_i(Y')/\mu_i(R)$。置 $X_i=\operatorname{coker}(p_1,p_2)$、$R'=Y'\times_{X'_i}Y'$。在 $F_i$ 中有\Term{canonical}{典范}\Term{morphism}{态射} $u:R\to R'$，且 $\mu_i^*(u)$ 是\Term{isomorphism}{同构}。因此改变指标 $i$ 后可假设 $R=R'$，即 $R$ 是\Term{equivalence-relation}{等价关系}。此时 $X_i$ 属于 $F_i^{\mathrm{coh}}$（1.17.1），所以 $X=\mu_i(X_i)$ 属于 $\overrightarrow F_{\mathrm{coh}}$。
\end{proof}
\begin{SGAresult}{推论}{8.3.14}
设 $I$ 是\Term{essentially-small}{本质小}的\Term{cofiltered-category}{余滤过范畴}，$p:F\to I$、$q:G\to I$ 是 $I$ 上的两个 $\mathcal U$-\Term{fibered-topos}{纤维拓扑斯}。假设每个 $i\in\operatorname{ob}I$ 处的 $F_i,G_i$ 都\Term{coherent-topos}{凝聚}，每个 $f:i\to j$ 的\Term{morphism}{态射} $f_\bullet:F_i\to F_j$ 和 $G_i\to G_j$ 都\Term{coherent-morphism}{凝聚}。又设 $m:F\to G$ 是 $\mathcal U$-\Term{fibered-topos}{纤维拓扑斯}的一个\Term{cartesian-morphism}{笛卡尔态射}，使每个 $m_i:F_i\to G_i$（$i\in\operatorname{ob}I$）都\Term{coherent-morphism}{凝聚}。则由 $m$ 取\Term{projective-limit}{投射极限}而得的\Term{morphism}{态射} $\underline m$ \Term{coherent-morphism}{凝聚}。
\end{SGAresult}
\begin{proof}
对每个\Term{object}{对象} $i$，$m_i$ 诱导\Term{functor}{函子} $m_{i,\mathrm{coh}}^*:G_{i,\mathrm{coh}}\to F_{i,\mathrm{coh}}$，从而给出\Term{cartesian-functor}{笛卡尔函子} $m_{\mathrm{coh}}^*:G_{\mathrm{coh}}\to F_{\mathrm{coh}}$；它是从\Term{fibered-site}{纤维景} $F_{\mathrm{coh}}$ 到 $G_{\mathrm{coh}}$ 的一个\Term{morphism}{态射}（7.4.13）。\Term{functor}{函子}
\[\varinjlim_{I^\circ}m_{\mathrm{coh}}^*:\varinjlim_{I^\circ}G_{\mathrm{coh}}\longrightarrow\varinjlim_{I^\circ}F_{\mathrm{coh}}\]
是\Term{site-morphism}{景态射}，转到相应\Term{topos}{拓扑斯}便得到\Term{morphism}{态射} $\underline m$。断言由 8.3.13 及 3.3 得到。
\end{proof}
\begin{SGAresult}{练习}{8.3.15}
采用 8.3.13 的记号，证明若各 $F_i$（$i\in I$）是\Term{algebraic-topos}{代数拓扑斯}（2.3），且各\Term{morphism}{态射} $f_\bullet:F_i\to F_j$（$f\in\operatorname{Fl}(I)$）\Term{coherent-morphism}{凝聚}，则 $\varprojlim_I\operatorname{top}F_i$
\SourcePage{220}
是\Term{algebraic-topos}{代数拓扑斯}，且
\[\left(\varprojlim_I\!\operatorname{top}F_i\right)_{\mathrm{coh}}\xrightarrow\sim\overrightarrow F_{\mathrm{coh}}.\]
\end{SGAresult}
